\chapter{Microlocalization}\label{ch:microlocal}
\providecommand{\Tz}[1]{T^*#1\setminus0}% cotangent space minus the zero section
\providecommand{\conesupp}{\operatorname{cone\,supp}}% cone support of a symbol
\providecommand{\Ell}{\operatorname{Ell}}% elliptic set

The pseudodifferential operators of \cref{ch:euclidean} are not local: the
support of $Au$ is in general not contained in the support of $u$. They are,
however, \emph{pseudolocal}: they do not create new singularities. This
chapter refines that statement in two directions. First, singularities of a
distribution are not only located at points of $\R^n$ but also have
directions: the \emph{wavefront set} $\WF(u)$ is a closed conic subset of
$\Tz{\R^n}=\R^n\times(\R^n\setminus0)$ which projects onto the singular
support. Secondly, an operator $A\in\Psido m(\R^n)$ has an \emph{essential
support} $\WFp(A)\subset\Tz{\R^n}$, the set where its symbol is not rapidly
decreasing, and $A$ is \emph{microlocal}:
$\WF(Au)\subset\WFp(A)\cap\WF(u)$. Conversely, where $A$ is elliptic it has
a microlocal parametrix, so $\WF(u)\subset\WF(Au)\cup\Char_m(A)$.

\Cref{microlocal:sec:supports,microlocal:sec:pseudolocal} deal with supports
and singular supports, \cref{microlocal:sec:ess,microlocal:sec:char,microlocal:sec:parametrix}
with the essential support, ellipticity at a point and microlocal
parametrices, and \cref{microlocal:sec:coordinates} with the invariance of
the calculus under changes of coordinates; this is what makes it possible to
transfer everything to manifolds in \cref{ch:manifolds}. The wavefront set
is introduced in \cref{microlocal:sec:wf} and characterized in terms of the
Fourier transform in \cref{microlocal:sec:fourier}. The remaining sections
develop the calculus of wavefront sets: invariance under diffeomorphisms,
the wavefront set of the kernel of a pseudodifferential operator, and the
behaviour of wavefront sets under pairing, exterior products, restriction,
projection, multiplication, pull-back and push-forward, culminating in the
wavefront relation of an operator with a distributional kernel.

\subsection*{Conventions}
We use freely the results of \cref{ch:distributions,ch:euclidean}. We write
$u(\phi)$ for the (bilinear) pairing of a distribution with a test function
and $\hat u=\cF u$ for the Fourier transform. A function $f$ on a set
$U\subset\R^p\times\R^n$ is \emph{rapidly decreasing on $U$} if
$\sup_{(x,\xi)\in U}\ang\xi^N\abs{f(x,\xi)}<\infty$ for every $N$; similarly
for functions of $\xi$ alone. For $\bar\xi\in\R^n\setminus0$, $\bar x\in\R^p$
and $\eps>0$ we use the open cones
\begin{equation*}
  V_\eps(\bar\xi)=\Bigl\{\xi\in\R^n\setminus0;\
  \Bigl|\frac\xi{\abs\xi}-\frac{\bar\xi}{\abs{\bar\xi}}\Bigr|<\eps\Bigr\},\qquad
  \Gamma_\eps(\bar x,\bar\xi)=\{x;\abs{x-\bar x}<\eps\}\times V_\eps(\bar\xi).
\end{equation*}
A set $W\subset\R^p\times(\R^n\setminus0)$ is \emph{conic} if
$(x,\xi)\in W\Rightarrow(x,t\xi)\in W$ for all $t>0$. Every conic
neighbourhood of $(\bar x,\bar\xi)$ contains some $\Gamma_\eps(\bar x,\bar\xi)$.
For a conic set $W$ we put $W^-=\{(x,-\xi);(x,\xi)\in W\}$. We shall often use
the elementary inequality
\begin{equation}\label{microlocal:eq:angle}
  \abs{se-tf}\ge\tfrac12\abs{e-f}\,(s+t),\qquad s,t\ge0,\ \abs e=\abs f=1,
\end{equation}
which follows from
$\abs{se-tf}^2=(s-t)^2+st\abs{e-f}^2\ge\frac12\abs{e-f}^2(s^2+t^2)$ (note
$\abs{e-f}\le2$), and the inequality
$\bigl|\frac v{\abs v}-\frac w{\abs w}\bigr|\le\frac{2\abs{v-w}}{\abs w}$ for
$v,w\ne0$.

From \cref{ch:euclidean} we recall in particular: every $A\in\Psido m(\R^n)$
is $A=\Op(\sigma_L(A))=\Op_R(\sigma_R(A))$ with unique left and right
symbols in $S^m(\R^n;\R^n)$; if $A=\Op(a)$ for an amplitude
$a\in S^m(\R^{2n};\R^n)$ (an amplitude of weight $0$), then its kernel is the
oscillatory integral $K_A=I(a)=\int e^{i(x-y)\cdot\xi}a(x,y,\xi)\,\dbar\xi$
and
\begin{equation}\label{microlocal:eq:reduction}
  \sigma_L(A)(x,\xi)\sim\sum_\alpha\frac{i^{\abs\alpha}}{\alpha!}
  D^\alpha_\xi D^\alpha_ya(x,y,\xi)\big|_{y=x};
\end{equation}
$\Op(a)\Op_R(b)=\Op\bigl(a(x,\xi)b(y,\xi)\bigr)$ for symbols $a$, $b$;
$\Psido{-\infty}(\R^n)$ consists of the operators whose kernels are smooth
and satisfy $\abs{D^\alpha_xD^\beta_yK(x,y)}\le C_{\alpha\beta N}\ang{x-y}^{-N}$
for all $\alpha,\beta,N$; such operators map $\cSp(\R^n)$ into $\CI(\R^n)$,
with $Eu(x)=u(K_E(x,\cdot))$; every $A\in\Psido m(\R^n)$ acts continuously
on $\cS(\R^n)$ and on $\cSp(\R^n)$, the latter by
$(Au)(\phi)=u(A^t\phi)$; and an operator $A\in\Psido m(\R^n)$ which is
\emph{(globally) elliptic}, in the sense that
$\abs{\sigma_L(A)(x,\xi)}\ge\eps\ang\xi^m$ for $\abs\xi\ge R$ and some
$\eps,R>0$, has a parametrix $B\in\Psido{-m}(\R^n)$ with
$AB-\Id,\ BA-\Id\in\Psido{-\infty}(\R^n)$.

% ======================================================================
\section{Supports and singular supports}\label{microlocal:sec:supports}

Recall from \cref{ch:distributions} that the support $\supp u$ of
$u\in\cSp(\R^n)$ is the complement of the largest open set on which $u$
vanishes; equivalently, $\bar x\notin\supp u$ if and only if there exists
$\phi\in\CIc(\R^n)$ with $\phi(\bar x)\ne0$ and $\phi u=0$. We write
\begin{equation*}
  \CmIc(\R^n)=\{u\in\cSp(\R^n);\ \supp u\text{ compact}\}.
\end{equation*}
An element $u\in\CmIc(\R^n)$ acts on all of $\CI(\R^n)$ by
$u(\phi)=u(\psi\phi)$, where $\psi\in\CIc(\R^n)$ equals $1$ near $\supp u$;
this does not depend on $\psi$. Since $u$ is tempered,
$\abs{u(\phi)}\le C\norm{\psi\phi}_k$, so $u$ has \emph{finite order}: for a
compact neighbourhood $L$ of $\supp u$,
\begin{equation}\label{microlocal:eq:finite-order}
  \abs{u(\phi)}\le C\sum_{\abs\beta\le k}\sup_{L}\abs{D^\beta\phi},
  \qquad\phi\in\CI(\R^n).
\end{equation}

\subsection{The calculus of supports}
If $A:\cS(\R^n)\to\cSp(\R^n)$ is continuous and linear, its Schwartz kernel
$K_A\in\cSp(\R^{2n})$ has a support, and we set $\supp(A)=\supp(K_A)$. A
subset $G\subset\R^n\times\R^n$ defines a \emph{relation}, that is a map on
subsets,
\begin{equation*}
  G\circ S=\{x\in\R^n;\ \exists\,y\in S\text{ with }(x,y)\in G\},
  \qquad S\subset\R^n.
\end{equation*}
In terms of the projections $\pi_L(x,y)=x$, $\pi_R(x,y)=y$,
\begin{equation*}
  \begin{tikzcd}[column sep=small]
    & \R^{2n}\arrow[dl,"\pi_L"']\arrow[dr,"\pi_R"] & \\
    \R^n & & \R^n
  \end{tikzcd}
  \qquad
  G\circ S=\pi_L\bigl(\pi_R^{-1}(S)\cap G\bigr).
\end{equation*}
If $S$ is compact and $G$ is closed then $G\circ S$ is closed: the restriction
of $\pi_L$ to the closed set $\pi_R^{-1}(S)\cap G$ is proper (the preimage of
a compact set $K$ is a closed subset of $K\times S$), and proper continuous
maps into $\R^n$ are closed.

\begin{proposition}[Calculus of supports]\label{microlocal:prop:supports}
  If $A:\cS(\R^n)\to\cSp(\R^n)$ is continuous and linear then
  \begin{equation*}
    \supp(A\phi)\subset\supp(A)\circ\supp(\phi)\quad\text{for all }\phi\in\CIc(\R^n).
  \end{equation*}
\end{proposition}
\begin{proof}
  Since $\supp(A)\circ\supp(\phi)$ is closed, a point $\bar x$ outside it has
  a neighbourhood disjoint from it; choose $\psi\in\CIc(\R^n)$ supported in
  this neighbourhood with $\psi(\bar x)\ne0$. If $(x,y)\in\supp(A)$ with
  $y\in\supp\phi$ then $x\in\supp(A)\circ\supp(\phi)$, so $x\notin\supp\psi$.
  Thus $\supp(A)\cap(\supp\psi\times\supp\phi)=\emptyset$. For every
  $\chi\in\cS(\R^n)$ the function $(\chi\psi)\boxtimes\phi$ vanishes near
  $\supp K_A$, hence
  $(\psi A\phi)(\chi)=K_A\bigl((\chi\psi)\boxtimes\phi\bigr)=0$. So
  $\psi A\phi=0$ and $\bar x\notin\supp(A\phi)$.
\end{proof}

For instance, the kernel of a differential operator is supported in the
diagonal $\Diag=\{(x,x)\}$, so differential operators are \emph{local}:
$\supp(P\phi)\subset\supp\phi$. A pseudodifferential operator, on the
other hand, generally has $\supp(A)=\R^{2n}$; the kernel of $\ang D^{-2}$ on
$\R$, for example, is $\frac12e^{-\abs{x-y}}$.

\subsection{Singular support}

\begin{definition}\label{microlocal:def:singsupp}
  The \emph{singular support} of $u\in\cSp(\R^n)$ is the set $\singsupp u$
  of points $\bar x$ such that there is no $\phi\in\CIc(\R^n)$ with
  $\phi(\bar x)\ne0$ and $\phi u\in\CI(\R^n)$.
\end{definition}

If $\phi u\in\CI$ and $\phi(\bar x)\ne0$ then $\phi\ne0$ on a ball $B$
around $\bar x$, and for every $\psi\in\CIc(B)$,
$\psi u=(\psi/\phi)\phi u\in\CI$. Hence the complement of $\singsupp u$ is
open, $\singsupp u$ is closed, and $\bar x\notin\singsupp u$ if and only if
$\phi u\in\CIc(\R^n)$ for some $\phi\in\CIc(\R^n)$ with $\phi=1$ near $\bar
x$. Directly from the definitions,
\begin{equation*}
  \singsupp u\subset\supp u,\qquad
  \singsupp u=\emptyset\iff u\in\CI(\R^n),
\end{equation*}
the second because $u$ is smooth if it is smooth near each point.

% ======================================================================
\section{Pseudolocality}\label{microlocal:sec:pseudolocal}

The first observation is that the kernel of a pseudodifferential operator is
smooth, and rapidly decreasing, away from the diagonal.

\begin{lemma}\label{microlocal:lem:off-diagonal}
  Let $a\in S^m(\R^{2n};\R^n)$ be an amplitude and let $\rho\in\CI(\R^n)$
  have all derivatives bounded and vanish near $0$. Then
  $\rho(x-y)I(a)(x,y)$ is the kernel of an element of
  $\Psido{-\infty}(\R^n)$: it is smooth and
  \begin{equation*}
    \abs{D^\beta_xD^\gamma_y\bigl(\rho(x-y)I(a)\bigr)}
    \le C_{\beta\gamma N}\ang{x-y}^{-N}\quad\text{for all }\beta,\gamma,N.
  \end{equation*}
\end{lemma}
\begin{proof}
  By the kernel identities of \cref{ch:euclidean},
  \begin{equation}\label{microlocal:eq:kernel-ids}
    D^\beta_xD^\gamma_y\bigl[(x-y)^\alpha I(a)\bigr]
    =I\bigl((\xi+D_x)^\beta(D_y-\xi)^\gamma(-D_\xi)^\alpha a\bigr),
  \end{equation}
  and the amplitude on the right has order $m+\abs\beta+\abs\gamma-\abs\alpha$.
  If this is less than $-n$ the oscillatory integral converges absolutely
  and uniformly, so for $\abs\alpha=N>m+n+k$ the function
  $k_\alpha=(x-y)^\alpha I(a)$ is of class $\mathcal C^k$ with all
  derivatives of order at most $k$ bounded.

  Suppose $\rho(z)=0$ for $\abs z<\delta$. Using
  $\abs z^{2N}=\sum_{\abs\alpha=N}\frac{N!}{\alpha!}z^{2\alpha}$ we can write
  \begin{equation*}
    \rho(z)=\sum_{\abs\alpha=N}\frac{N!}{\alpha!}h_\alpha(z)z^\alpha,\qquad
    h_\alpha(z)=\frac{\rho(z)z^\alpha}{\abs z^{2N}},
  \end{equation*}
  where $h_\alpha\in\CI(\R^n)$ and, since $z^\alpha\abs z^{-2N}$ is
  homogeneous of degree $-N$ on $\abs z\ge\delta$,
  $\abs{D^\mu h_\alpha(z)}\le C_\mu\ang z^{-N}$. Then
  $\rho(x-y)I(a)=\sum_{\abs\alpha=N}\frac{N!}{\alpha!}h_\alpha(x-y)k_\alpha$ as
  distributions, and by Leibniz' formula this is a $\mathcal C^k$ function
  whose derivatives of order at most $k$ are bounded by $C\ang{x-y}^{-N}$.
  Since $k$ and $N>m+n+k$ are arbitrary, the lemma follows.
\end{proof}

\begin{proposition}\label{microlocal:prop:kernel-singsupp}
  If $A\in\Psido m(\R^n)$ then $\singsupp(K_A)\subset\Diag$. Moreover, if
  $\phi,\psi\in\CIb(\R^n)$ have supports at positive distance from each
  other, then $\phi A\psi\in\Psido{-\infty}(\R^n)$.
\end{proposition}
\begin{proof}
  Write $K_A=I(a)$ with $a=\sigma_L(A)(x,\xi)$. If $\bar x\ne\bar y$, choose
  $\rho$ as in \cref{microlocal:lem:off-diagonal} with $\rho=1$ near
  $\bar x-\bar y$; then $K_A=\rho(x-y)K_A$ near $(\bar x,\bar y)$, which is
  smooth. For the second statement, if $\abs{x-y}\ge\delta>0$ whenever
  $x\in\supp\phi$, $y\in\supp\psi$, choose $\rho$ with $\rho(z)=1$ for
  $\abs z\ge\delta$ and $\rho(z)=0$ for $\abs z\le\delta/2$. The kernel of
  $\phi A\psi$ is
  $\phi(x)\psi(y)K_A=\phi(x)\psi(y)\rho(x-y)K_A$, and multiplication by
  $\phi(x)\psi(y)$, a function with bounded derivatives, preserves the
  estimates of \cref{microlocal:lem:off-diagonal}.
\end{proof}

We also note how residual operators act on distributions of compact
support.

\begin{lemma}\label{microlocal:lem:residual-compact}
  If $E\in\Psido{-\infty}(\R^n)$ and $u\in\CmIc(\R^n)$ then $Eu\in\cS(\R^n)$.
\end{lemma}
\begin{proof}
  $Eu(x)=u(K_E(x,\cdot))$, and differentiating under $u$,
  $x^\alpha D^\gamma Eu(x)=u\bigl(x^\alpha D^\gamma_xK_E(x,\cdot)\bigr)$. By
  \eqref{microlocal:eq:finite-order} this is bounded by
  $C\sup_{y\in L,\abs\beta\le k}\abs{x^\alpha D^\gamma_xD^\beta_yK_E(x,y)}$,
  which is finite since $\ang x\le C_L\ang{x-y}$ for $y$ in the compact set
  $L$ and $\abs{D^\gamma_xD^\beta_yK_E}\le C_N\ang{x-y}^{-N}$.
\end{proof}

\begin{definition}
  A continuous linear operator $A:\cSp(\R^n)\to\cSp(\R^n)$ is
  \emph{pseudolocal} if $\singsupp(Au)\subset\singsupp(u)$ for all
  $u\in\cSp(\R^n)$.
\end{definition}

\begin{proposition}\label{microlocal:prop:pseudolocal}
  Pseudodifferential operators are pseudolocal.
\end{proposition}
\begin{proof}
  Let $A\in\Psido m(\R^n)$, $u\in\cSp(\R^n)$ and $\bar x\notin\singsupp u$.
  Choose $\phi\in\CIc(\R^n)$ with $\phi=1$ on a neighbourhood $U$ of $\bar x$
  and $\phi u\in\CIc(\R^n)$, and $\psi\in\CIc(U)$ with $\psi=1$ near $\bar x$.
  Then
  \begin{equation*}
    \psi Au=\psi A(\phi u)+\psi A(1-\phi)u.
  \end{equation*}
  The first term is in $\cS(\R^n)$, since $A$ preserves $\cS$. Since
  $\supp\psi$ is a compact subset of the open set $U$ and $1-\phi=0$ on $U$,
  the supports of $\psi$ and $1-\phi$ are at positive distance, so
  $\psi A(1-\phi)\in\Psido{-\infty}(\R^n)$ by
  \cref{microlocal:prop:kernel-singsupp}, and the second term is smooth.
  Hence $\psi Au\in\CI$ and $\bar x\notin\singsupp(Au)$.
\end{proof}

% ======================================================================
\section{Cone support and essential support}\label{microlocal:sec:ess}

\begin{definition}\label{microlocal:def:conesupp}
  The \emph{cone support} of $a\in S^m(\R^p;\R^n)$ is the set
  $\conesupp(a)\subset\R^p\times(\R^n\setminus0)$ of points $(\bar x,\bar\xi)$
  such that $a$ is \emph{not} rapidly decreasing on any set
  $\Gamma_\eps(\bar x,\bar\xi)$, $\eps>0$.
\end{definition}

The complement of $\conesupp(a)$ is open and conic, so $\conesupp(a)$ is a
closed conic subset of $\R^p\times(\R^n\setminus0)$. Only the behaviour of
$a$ for large $\abs\xi$ matters: changing $a$ on a set where $\abs\xi$ is
bounded does not change the cone support. A priori the definition only
involves $a$ itself, but in fact all derivatives decay as well.

\begin{lemma}\label{microlocal:lem:conesupp-derivatives}
  Let $a\in S^m(\R^p;\R^n)$ be rapidly decreasing on $\Gamma_\eps(\bar
  x,\bar\xi)$. Then for every $\eps'<\eps$ and all $\alpha,\beta$,
  $D^\alpha_xD^\beta_\xi a$ is rapidly decreasing on $\Gamma_{\eps'}(\bar
  x,\bar\xi)$. In particular
  $\conesupp(D^\alpha_xD^\beta_\xi a)\subset\conesupp(a)$.
\end{lemma}
\begin{proof}
  Since $D_{x_j}a\in S^m$ and $D_{\xi_j}a\in S^{m-1}$, by induction (with
  finitely many intermediate values of $\eps$ between $\eps'$ and $\eps$) it
  suffices to show that $\pa_{x_j}a$ and $\pa_{\xi_j}a$ are rapidly
  decreasing on $\Gamma_{\eps''}=\Gamma_{\eps''}(\bar x,\bar\xi)$ for any
  $\eps''<\eps$. Fix $M\ge\abs m$; we show that both derivatives are
  $O(\ang\xi^{(m-M)/2})$ on $\Gamma_{\eps''}$, which suffices as $M$ is
  arbitrary. Since all derivatives of $a$ are $O(\ang\xi^m)$ we may restrict
  to $\abs\xi\ge R$, for an $R$ to be chosen.

  Let $(x,\xi)\in\Gamma_{\eps''}$ and $0<t\le\eps-\eps''$. Then
  $(x+se_j,\xi)\in\Gamma_\eps$ for $\abs s\le t$, and Taylor's formula gives
  \begin{equation*}
    \abs{t\,\pa_{x_j}a(x,\xi)}\le\abs{a(x+te_j,\xi)}+\abs{a(x,\xi)}
    +\tfrac12t^2\sup_{\abs s\le t}\abs{\pa^2_{x_j}a(x+se_j,\xi)}
    \le 2C_M\ang\xi^{-M}+Ct^2\ang\xi^m.
  \end{equation*}
  Taking $t=(\eps-\eps'')\ang\xi^{-(M+m)/2}$ gives
  $\abs{\pa_{x_j}a(x,\xi)}\le C'\ang\xi^{(m-M)/2}$.

  For the $\xi$-derivative let $0<t\le1$ and $\abs s\le t$. Then
  $\bigl|\frac{\xi+se_j}{\abs{\xi+se_j}}-\frac\xi{\abs\xi}\bigr|\le\frac2{\abs\xi}$,
  so $(x,\xi+se_j)\in\Gamma_\eps$ if $\abs\xi\ge R=2/(\eps-\eps'')+1$; also
  $\ang{\xi+se_j}\ge\frac12\ang\xi$. As before
  \begin{equation*}
    \abs{t\,\pa_{\xi_j}a(x,\xi)}\le C\ang\xi^{-M}+Ct^2\ang\xi^{m},
  \end{equation*}
  and $t=\ang\xi^{-(M+m)/2}$ gives $\abs{\pa_{\xi_j}a}\le C'\ang\xi^{(m-M)/2}$.
\end{proof}

\begin{corollary}\label{microlocal:cor:conesupp-compact}
  Let $a\in S^m(\R^p;\R^n)$ and let $K\subset\R^p$ be compact with
  $\conesupp(a)\cap(K\times(\R^n\setminus0))=\emptyset$. Then $\phi a\in
  S^{-\infty}(\R^p;\R^n)$ for every $\phi\in\CIc(\R^p)$ with $\supp\phi\subset
  K$. In particular, if $a$ has compact support in $x$ and
  $\conesupp(a)=\emptyset$, then $a\in S^{-\infty}(\R^p;\R^n)$.
\end{corollary}
\begin{proof}
  Each point of the compact set $K\times\Sphere^{n-1}$ has a neighbourhood
  $\Gamma_{\eps}(\bar x,\bar\xi)$ on which $a$ is rapidly decreasing, hence
  (by \cref{microlocal:lem:conesupp-derivatives}) with all derivatives on
  $\Gamma_{\eps/2}(\bar x,\bar\xi)$. Finitely many of the latter cover
  $K\times(\R^n\setminus0)$, so every $D^\alpha_xD^\beta_\xi a$ is rapidly
  decreasing on $K\times\R^n$, and by Leibniz' formula $\phi a\in
  S^{-\infty}$.
\end{proof}

\begin{remark}\label{microlocal:rem:conesupp-global}
  The compactness in $x$ cannot be dropped: a symbol may have empty cone
  support without being in $S^{-\infty}$, because the definition only
  controls it locally uniformly in $x$; see \cref{microlocal:ex:empty-wf}.
\end{remark}

\begin{definition}\label{microlocal:def:ess}
  The \emph{essential support} (or \emph{operator wavefront set}) of
  $A\in\Psido m(\R^n)$ is the closed conic set
  \begin{equation*}
    \WFp(A)=\conesupp\bigl(\sigma_L(A)\bigr)\subset\Tz{\R^n}.
  \end{equation*}
\end{definition}

The essential support can be computed from any amplitude.

\begin{lemma}\label{microlocal:lem:ess-amplitude}
  If $A=\Op(a)$ with $a\in S^m(\R^{2n};\R^n)$ then
  \begin{equation*}
    \WFp(A)\subset\{(x,\xi);\ (x,x,\xi)\in\conesupp(a)\}.
  \end{equation*}
\end{lemma}
\begin{proof}
  Suppose $(\bar x,\bar x,\bar\xi)\notin\conesupp(a)$, so $a$ is rapidly
  decreasing on $\{\abs{x-\bar x}<\eps,\abs{y-\bar x}<\eps,\xi\in
  V_\eps(\bar\xi)\}$ for some $\eps>0$. By
  \cref{microlocal:lem:conesupp-derivatives} (applied with $(x,y)$ as base
  variables) all derivatives of $a$ are rapidly decreasing on the
  corresponding set with $\eps/2$ in place of $\eps$, so each term
  $D^\alpha_\xi D^\alpha_ya(x,y,\xi)|_{y=x}$ in \eqref{microlocal:eq:reduction}
  is rapidly decreasing on $\Gamma_{\eps/2}(\bar x,\bar\xi)$. By
  \eqref{microlocal:eq:reduction},
  $\sigma_L(A)-\sum_{\abs\alpha<N}\frac{i^{\abs\alpha}}{\alpha!}D^\alpha_\xi
  D^\alpha_ya|_{y=x}\in S^{m-N}$, so
  $\abs{\sigma_L(A)}\le C_N\ang\xi^{m-N}$ on $\Gamma_{\eps/2}(\bar
  x,\bar\xi)$ for every $N$, that is $(\bar x,\bar\xi)\notin\WFp(A)$.
\end{proof}

\begin{proposition}\label{microlocal:prop:ess}
  For $A\in\Psido m(\R^n)$ and $B\in\Psido{m'}(\R^n)$:
  \begin{enumerate}[label=\textup{(\roman*)}]
  \item $\WFp(A)=\conesupp(\sigma_R(A))$;
  \item $\WFp(A+B)\subset\WFp(A)\cup\WFp(B)$ (if $m=m'$) and
    $\WFp(AB)\subset\WFp(A)\cap\WFp(B)$;
  \item $\WFp(A^*)=\WFp(A)$ and $\WFp(A^t)=\WFp(A)^-$;
  \item $\WFp(\Id)=\Tz{\R^n}$, $\WFp(E)=\emptyset$ for
    $E\in\Psido{-\infty}(\R^n)$, and for multiplication by $\phi\in\CIb(\R^n)$,
    $\WFp(\phi)=\supp\phi\times(\R^n\setminus0)$.
  \end{enumerate}
\end{proposition}
\begin{proof}
  (i) $A=\Op(b(y,\xi))$ with $b=\sigma_R(A)$, and $(x,x,\xi)\in\conesupp(b(y,\xi))$
  if and only if $(x,\xi)\in\conesupp(b)$, so $\WFp(A)\subset\conesupp(\sigma_R(A))$
  by \cref{microlocal:lem:ess-amplitude}. Conversely
  $\sigma_R(A)\sim\sum_\beta\frac{(-i)^{\abs\beta}}{\beta!}D^\beta_\xi
  D^\beta_x\sigma_L(A)$, and the argument in the proof of
  \cref{microlocal:lem:ess-amplitude} shows
  $\conesupp(\sigma_R(A))\subset\conesupp(\sigma_L(A))$.

  (ii) The first statement is clear. For the second, let $a=\sigma_L(A)$ and
  $b=\sigma_R(B)$, so that $AB=\Op(a(x,\xi)b(y,\xi))$. If $(\bar
  x,\bar\xi)\notin\WFp(A)$ then $a(x,\xi)b(y,\xi)$ is rapidly decreasing on
  $\{\abs{x-\bar x}<\eps,\xi\in V_\eps(\bar\xi)\}$ for some $\eps$, since $b$
  is polynomially bounded; if $(\bar x,\bar\xi)\notin\WFp(B)=\conesupp(b)$,
  the same holds with the roles of $x$ and $y$ exchanged. In either case
  $(\bar x,\bar x,\bar\xi)\notin\conesupp(a(x,\xi)b(y,\xi))$, and
  \cref{microlocal:lem:ess-amplitude} applies.

  (iii) The adjoint of $\Op(a)$, $a=\sigma_L(A)$, is $\Op(\overline{a(y,\xi)})$,
  and its transpose is $\Op(a(y,-\xi))$ (the kernel of $A^t$ is
  $K_A(y,x)=\int e^{i(x-y)\cdot\xi}a(y,-\xi)\,\dbar\xi$). By
  \cref{microlocal:lem:ess-amplitude}, $\WFp(A^*)\subset\WFp(A)$ and
  $\WFp(A^t)\subset\WFp(A)^-$; applying this to $A^*$ and $A^t$ gives
  equality.

  (iv) is immediate from the definition, since $\sigma_L(\Id)=1$,
  $\sigma_L(E)\in S^{-\infty}$ and $\sigma_L(\phi)=\phi(x)$.
\end{proof}

\begin{lemma}\label{microlocal:lem:localize}
  Let $A\in\Psido m(\R^n)$ and $\phi\in\CIc(\R^n)$ with
  $\WFp(A)\cap(\supp\phi\times(\R^n\setminus0))=\emptyset$. Then $\phi A$ and
  $A\phi$ lie in $\Psido{-\infty}(\R^n)$. In particular, if $\WFp(A)=\emptyset$
  then $\phi A,A\phi\in\Psido{-\infty}(\R^n)$ for all $\phi\in\CIc(\R^n)$.
\end{lemma}
\begin{proof}
  $\sigma_L(\phi A)=\phi\sigma_L(A)\in S^{-\infty}$ by
  \cref{microlocal:cor:conesupp-compact}. Similarly
  $A\phi=\Op\bigl(\sigma_R(A)(y,\xi)\phi(y)\bigr)=\Op_R(\phi\sigma_R(A))$ and
  $\phi\sigma_R(A)\in S^{-\infty}$, by
  \cref{microlocal:prop:ess}(i) and \cref{microlocal:cor:conesupp-compact}.
\end{proof}

Note that $\WFp(A)=\emptyset$ does \emph{not} imply
$A\in\Psido{-\infty}(\R^n)$ (\cref{microlocal:rem:conesupp-global}); this is
why compactly supported cut-offs appear in many arguments below.

\subsection{Microlocal cut-offs}
We shall constantly use operators which are ``microlocally the identity''
near a given closed cone and vanish microlocally away from it. Fix
$\chi\in\CIc(\R^n)$ with $0\le\chi\le1$, $\chi(t)=1$ for $\abs t\le\frac12$ and
$\chi(t)=0$ for $\abs t\ge1$. For $z=(\bar x,\bar\xi)\in\Tz{\R^n}$ and
$0<\eps<1$ set
\begin{equation}\label{microlocal:eq:gamma}
  \gamma_{z,\eps}(x,\xi)=\chi\Bigl(\frac{x-\bar x}\eps\Bigr)
  \chi\Bigl(\frac1\eps\Bigl(\frac\xi{\abs\xi}-\frac{\bar\xi}{\abs{\bar\xi}}\Bigr)\Bigr)
  \bigl(1-\chi(\eps\xi)\bigr).
\end{equation}
The last factor vanishes for $\abs\xi\le1/(2\eps)$, and functions which are
smooth and homogeneous of degree $0$ in $\abs\xi\ge c>0$ are symbols of
order $0$ there, so $\gamma_{z,\eps}\in S^0(\R^n;\R^n)$. It takes values in
$[0,1]$, is supported in $\overline{\Gamma_\eps(z)}$ and equals $1$ on
$\Gamma_{\eps/2}(z)\cap\{\abs\xi\ge1/\eps\}$. Hence
$G_{z,\eps}=\Op(\gamma_{z,\eps})\in\Psido0(\R^n)$ satisfies
\begin{equation}\label{microlocal:eq:G-cutoff}
  \WFp(G_{z,\eps})\subset\overline{\Gamma_\eps(z)},\qquad
  \WFp(\Id-G_{z,\eps})\cap\Gamma_{\eps/2}(z)=\emptyset,\qquad
  G_{z,\eps}=\psi G_{z,\eps}
\end{equation}
for any $\psi\in\CIc(\R^n)$ equal to $1$ on $\{\abs{x-\bar x}\le\eps\}$.

\begin{lemma}[Microlocal cut-off]\label{microlocal:lem:cutoff}
  Let $W_1,W_2\subset\Tz{\R^n}$ be disjoint closed conic sets with $\pi(W_1)$
  bounded, where $\pi(x,\xi)=x$. Then there is $a\in S^0(\R^n;\R^n)$,
  real-valued with compact support in $x$, such that $C=\Op(a)$ satisfies
  \begin{equation*}
    \WFp(C)\cap W_2=\emptyset,\qquad\WFp(\Id-C)\cap W_1=\emptyset .
  \end{equation*}
\end{lemma}
\begin{proof}
  The set $W_1\cap(\R^n\times\Sphere^{n-1})$ is compact. For each of its
  points $z$ there is $\eps_z>0$ with $\overline{\Gamma_{2\eps_z}(z)}\cap
  W_2=\emptyset$, since $W_2$ is closed and conic. Choose finitely many
  $z_1,\dots,z_N$ such that the cones $\Gamma_{\eps_j/2}(z_j)$,
  $\eps_j=2\eps_{z_j}$, cover $W_1$, and put
  $g=\sum_j\gamma_{z_j,\eps_j}\in S^0$. Then $g\ge0$, $g\ge1$ on
  $\bigcup_j\Gamma_{\eps_j/2}(z_j)\cap\{\abs\xi\ge R\}$ with
  $R=\max_j1/\eps_j$, and $\supp g\subset\bigcup_j\overline{\Gamma_{\eps_j}(z_j)}$.
  Let $\Phi\in\CI(\R)$ with $\Phi(t)=0$ for $t\le0$ and $\Phi(t)=1$ for
  $t\ge1$, and set $a=\Phi\circ g$. By the chain rule $a\in S^0$, and
  $a$ has compact support in $x$. Since $a$ vanishes outside
  $\bigcup_j\overline{\Gamma_{\eps_j}(z_j)}$, which is a closed cone disjoint
  from $W_2$, $\WFp(C)\cap W_2=\emptyset$. Since $1-a$ vanishes on the open
  cone $\bigcup_j\Gamma_{\eps_j/2}(z_j)$ for $\abs\xi\ge R$, and this cone
  contains $W_1$, $\WFp(\Id-C)\cap W_1=\emptyset$.
\end{proof}

% ======================================================================
\section{The characteristic variety}\label{microlocal:sec:char}

\begin{definition}\label{microlocal:def:elliptic}
  $A\in\Psido m(\R^n)$ is \emph{elliptic at} $(\bar x,\bar\xi)\in\Tz{\R^n}$
  if there is $\eps>0$ such that
  \begin{equation*}
    \abs{\sigma_L(A)(x,\xi)}\ge\eps\abs\xi^m\quad\text{on }
    \Gamma_\eps(\bar x,\bar\xi)\cap\{\abs\xi\ge1/\eps\}.
  \end{equation*}
  The set of such points is the \emph{elliptic set} $\Ell_m(A)$; its
  complement $\Char_m(A)=\Tz{\R^n}\setminus\Ell_m(A)$ is the
  \emph{characteristic variety} of $A$.
\end{definition}

If $a'-\sigma_L(A)\in S^{m-1}$ then $\abs{a'-\sigma_L(A)}\le
C\ang\xi^{m-1}\le\frac\eps2\abs\xi^m$ for $\abs\xi$ large, so the condition
holds for $a'$ with $\eps/2$ in place of $\eps$ (after shrinking the cone).
Thus ellipticity at a point depends only on the principal symbol
$\sigma_m(A)$; in particular $\sigma_R(A)$ may be used instead of
$\sigma_L(A)$. Clearly $\Ell_m(A)$ is open and conic, so $\Char_m(A)$ is a
closed conic subset of $\Tz{\R^n}$. If $A$ is elliptic at $z$ then
$\sigma_L(A)$ is not rapidly decreasing near $z$, so
\begin{equation*}
  \Ell_m(A)\subset\WFp(A).
\end{equation*}
For example, multiplication by $\phi\in\CIb(\R^n)$ is elliptic (of order
$0$) at $(\bar x,\bar\xi)$ if and only if $\phi(\bar x)\ne0$, and
$G_{z,\eps}$ in \eqref{microlocal:eq:G-cutoff} is elliptic at $z$. The
operator $A$ is globally elliptic in the sense of \cref{ch:euclidean} if
and only if the estimate in \cref{microlocal:def:elliptic} holds uniformly,
with one $\eps$ for all $(\bar x,\bar\xi)$.

\begin{proposition}\label{microlocal:prop:char-product}
  For $A\in\Psido m(\R^n)$ and $B\in\Psido{m'}(\R^n)$,
  \begin{equation*}
    \Char_{m+m'}(AB)=\Char_m(A)\cup\Char_{m'}(B).
  \end{equation*}
\end{proposition}
\begin{proof}
  We may use the representatives $a=\sigma_L(A)$, $b=\sigma_L(B)$ and $ab$ of
  the principal symbols, since $\sigma_L(AB)-ab\in S^{m+m'-1}$. If $A$ and
  $B$ are elliptic at $z$ then $\abs{ab}\ge\eps^2\abs\xi^{m+m'}$ on a suitable
  $\Gamma_\eps(z)\cap\{\abs\xi\ge1/\eps\}$, so $AB$ is elliptic at $z$.
  Conversely if $\abs{ab}\ge\eps\abs\xi^{m+m'}$ there, then since
  $\abs b\le C\ang\xi^{m'}\le 2^{\abs{m'}}C\abs\xi^{m'}$ for $\abs\xi\ge1$,
  we get $\abs a\ge\eps'\abs\xi^m$, so $A$ is elliptic at $z$, and similarly
  $B$ is.
\end{proof}

% ======================================================================
\section{Microlocal parametrices}\label{microlocal:sec:parametrix}

The essential support allows us to refine the notion of a parametrix of an
elliptic operator to that of a parametrix at a point.

\begin{proposition}\label{microlocal:prop:parametrix}
  If $A\in\Psido m(\R^n)$ and $z\notin\Char_m(A)$, there exists a
  \emph{microlocal parametrix} for $A$ at $z$: an operator
  $B\in\Psido{-m}(\R^n)$ with
  \begin{equation*}
    z\notin\WFp(AB-\Id)\quad\text{and}\quad z\notin\WFp(BA-\Id).
  \end{equation*}
  Moreover there is then an open conic neighbourhood $\Gamma$ of $z$ with
  $\Gamma\cap\bigl(\WFp(AB-\Id)\cup\WFp(BA-\Id)\bigr)=\emptyset$.
\end{proposition}
\begin{proof}
  The last statement holds for any $B$ with the first property, as essential
  supports are closed. Let $a=\sigma_L(A)$ and choose $\eps\in(0,\frac12)$
  and $c_0>0$ such that $\abs a\ge c_0\abs\xi^m$ on
  $\Gamma_{2\eps}(z)\cap\{\abs\xi\ge1/(2\eps)\}$. Let $\gamma=\gamma_{z,\eps}$
  and $L=\Op(\gamma)$. Then $z\notin\WFp(\Id-L)$ by
  \eqref{microlocal:eq:G-cutoff}, and $\supp\gamma\subset
  \overline{\Gamma_\eps(z)}\cap\{\abs\xi\ge1/(2\eps)\}$, so that
  $\abs a\ge c_0\abs\xi^m$ on $\supp\gamma$.

  Let $\Lambda=\Op(\ang\xi^{2m})\in\Psido{2m}(\R^n)$ and consider
  \begin{equation*}
    J=(\Id-L)\Lambda+A^*A\in\Psido{2m}(\R^n).
  \end{equation*}
  Its principal symbol is represented by
  $(1-\gamma)\ang\xi^{2m}+\abs a^2$. Since $0\le\gamma\le1$, at a point
  with $\abs\xi\ge R$ either $\gamma\le\frac12$, and then this is at least
  $\frac12\ang\xi^{2m}$, or $\gamma>\frac12$, and then the point is in
  $\supp\gamma$ and the expression is at least
  $c_0^2\abs\xi^{2m}\ge c\ang\xi^{2m}$. Since $\sigma_L(J)$ differs from this
  representative by a symbol of order $2m-1$, $J$ is globally elliptic and
  has a parametrix $H\in\Psido{-2m}(\R^n)$, $HJ-\Id=E_1\in\Psido{-\infty}$.
  Set $B_1=HA^*\in\Psido{-m}(\R^n)$. Then
  \begin{equation*}
    B_1A-\Id=H\bigl(J-(\Id-L)\Lambda\bigr)-\Id=E_1-H(\Id-L)\Lambda,
  \end{equation*}
  and $\WFp(H(\Id-L)\Lambda)\subset\WFp(\Id-L)\not\ni z$ by
  \cref{microlocal:prop:ess}. So $z\notin\WFp(B_1A-\Id)$: $B_1$ is a left
  microlocal parametrix.

  In the same way $J'=(\Id-L)\Lambda+AA^*$ is globally elliptic; if $H'$ is
  a parametrix, $JH'-\Id=E_2\in\Psido{-\infty}$, then $B_2=A^*H'$ satisfies
  $AB_2-\Id=E_2-(\Id-L)\Lambda H'$, so $z\notin\WFp(AB_2-\Id)$. Write
  $R_1=B_1A-\Id$ and $R_2=AB_2-\Id$. Then
  \begin{equation*}
    B_2=(B_1A-R_1)B_2=B_1(\Id+R_2)-R_1B_2=B_1+B_1R_2-R_1B_2,
  \end{equation*}
  so $z\notin\WFp(B_2-B_1)$, and $AB_1-\Id=R_2+A(B_1-B_2)$ shows
  $z\notin\WFp(AB_1-\Id)$. Thus $B=B_1$ is a two-sided microlocal parametrix.
\end{proof}

The same argument shows that a microlocal parametrix at $z$ is unique up to
an operator whose essential support does not contain $z$
(\cref{microlocal:ex:uniqueness}).

% ======================================================================
\section{Coordinate invariance}\label{microlocal:sec:coordinates}

For an open set $\Omega\subset\R^n$ let $\CIc(\Omega)$ be the smooth functions
with compact support in $\Omega$ and
\begin{equation*}
  \CmIc(\Omega)=\{u\in\cSp(\R^n);\ \supp u\text{ is a compact subset of }\Omega\}.
\end{equation*}
An element $u\in\CmIc(\Omega)$ acts on $\CI(\Omega)$ by $u(\phi)=u(\psi\phi)$,
$\psi\in\CIc(\Omega)$ equal to $1$ near $\supp u$.

Let $F:\Omega\to\Omega'$ be a diffeomorphism between open subsets of $\R^n$,
with inverse $G$, and let $F^*\phi=\phi\circ F$, so
$F^*:\CIc(\Omega')\to\CIc(\Omega)$. If $u\in\CIc(\Omega')$ and
$\phi\in\CIc(\Omega)$, the change of variables $y=F(x)$ gives
\begin{equation}\label{microlocal:eq:pullback-change}
  \int u(F(x))\phi(x)\,dx=\int u(y)\,\phi(G(y))\,\abs{\det DG(y)}\,dy .
\end{equation}

\begin{lemma}\label{microlocal:lem:diffeo-pullback}
  For $u\in\CmIc(\Omega')$ the formula
  \begin{equation*}
    (F^*u)(\phi)=u\bigl(\abs{\det DG}\,G^*\phi\bigr),\qquad\phi\in\cS(\R^n),
  \end{equation*}
  (where $G^*\phi=\phi\circ G\in\CI(\Omega')$) defines $F^*u\in\CmIc(\Omega)$
  with $\supp F^*u=G(\supp u)$. It extends $F^*$ on $\CIc(\Omega')$, is
  continuous in the sense that $F^*u_j\to F^*u$ weakly if $u_j\to u$ weakly
  and the supports of the $u_j$ lie in a fixed compact subset of $\Omega'$,
  and satisfies $(F_1\circ F_2)^*=F_2^*F_1^*$; in particular
  $G^*F^*=\Id$.
\end{lemma}
\begin{proof}
  $u$ acts on $\CI(\Omega')$, and $\phi\mapsto\abs{\det DG}G^*\phi$ is
  continuous from $\cS(\R^n)$ to $\CI(\Omega')$, so $F^*u\in\cSp(\R^n)$. If
  $\phi$ vanishes near $G(\supp u)$ then $\abs{\det DG}G^*\phi$ vanishes near
  $\supp u$, so $\supp F^*u\subset G(\supp u)$, and equality follows from
  $G^*F^*u=u$. The consistency with smooth functions is
  \eqref{microlocal:eq:pullback-change}, continuity is clear, and the
  composition law follows from the chain rule for Jacobians.
\end{proof}

If $u\in\CmIc(\Omega')$ and $\rho_\eps(x)=\eps^{-n}\rho(x/\eps)$ with
$\rho\in\CIc(\R^n)$, $\int\rho=1$, then $\rho_\eps*u\in\CIc(\Omega')$ for small
$\eps$, with supports in a fixed compact set, and $\rho_\eps*u\to u$ weakly.
So $\CIc(\Omega')$ is weakly dense in $\CmIc(\Omega')$ in this strong sense
(see also \cref{microlocal:ex:density}), and $F^*$ is the unique such
continuous extension.

Now let $A:\cS(\R^n)\to\cSp(\R^n)$ have kernel with compact support in
$\Omega\times\Omega$. By \cref{microlocal:prop:supports}, $A$ maps
$\CIc(\Omega)$, indeed $\cS(\R^n)$, into $\CmIc(\Omega)$, and we define
\begin{equation}\label{microlocal:eq:AF}
  A_F=G^*\circ A\circ F^*:\CIc(\Omega')\to\CmIc(\Omega').
\end{equation}
By \eqref{microlocal:eq:pullback-change} its kernel is
\begin{equation*}
  K_{A_F}(x',y')=K_A(G(x'),G(y'))\,\abs{\det DG(y')},
\end{equation*}
that is, $K_{A_F}=\abs{\det DG(y')}\,(G\times G)^*K_A$, where $(G\times G)^*$
is the pull-back of \cref{microlocal:lem:diffeo-pullback} for the
diffeomorphism $G\times G:\Omega'\times\Omega'\to\Omega\times\Omega$; it has compact support in $\Omega'\times\Omega'$, and we extend it
by zero to $\R^{2n}$, so $A_F:\cS(\R^n)\to\cSp(\R^n)$. The diffeomorphism
$F$ induces the map
\begin{equation*}
  \chi_F:\Omega\times\R^n\to\Omega'\times\R^n,\qquad
  \chi_F(x,\xi)=\bigl(F(x),(DF(x)^t)^{-1}\xi\bigr),
\end{equation*}
which is a diffeomorphism, linear and invertible on each fibre. (In the
language of \cref{ch:manifolds}, $\chi_F$ is the inverse of the cotangent map
$F^*:T^*_{F(x)}\Omega'\to T^*_x\Omega$, $\eta\mapsto DF(x)^t\eta$.)

\begin{theorem}[Coordinate invariance]\label{microlocal:thm:coordinates}
  Let $F:\Omega\to\Omega'$ be a diffeomorphism of open subsets of $\R^n$ and
  $A\in\Psido m(\R^n)$ an operator whose kernel has compact support in
  $\Omega\times\Omega$. Then $A_F\in\Psido m(\R^n)$, its kernel has compact
  support in $\Omega'\times\Omega'$, and
  \begin{gather*}
    \sigma_m(A_F)\circ\chi_F=\sigma_m(A),\qquad
    \WFp(A_F)=\chi_F\bigl(\WFp(A)\bigr),\\
    \Char_m(A_F)\cap\pi^{-1}(\Omega')=\chi_F\bigl(\Char_m(A)\cap\pi^{-1}(\Omega)\bigr),
  \end{gather*}
  where $\pi(x,\xi)=x$ and the first identity is understood as
  $\sigma_L(A_F)(F(x),(DF(x)^t)^{-1}\xi)-\sigma_L(A)(x,\xi)\in S^{m-1}$ for
  $x$ in a neighbourhood of the support, both symbols vanishing for $x$
  outside compact subsets of $\Omega$, $\Omega'$ respectively.
\end{theorem}
\begin{proof}
  Let $K\subset\Omega$ be compact with $\supp K_A\subset K\times K$, and let
  $\psi\in\CIc(\Omega)$ equal $1$ near $K$. Then $\psi A=A$, so
  $a=\sigma_L(A)$ satisfies $a=\psi a$ by uniqueness of the left symbol, and
  $A=\psi A\psi=\Op(b_0)$ with $b_0(x,y,\xi)=\psi(x)\psi(y)a(x,\xi)$.
  Similarly $\sigma_L(A_F)$ is supported in $F(K)\times\R^n$.

  Let $L=\supp\psi$. Choose $\delta'>0$ so small that the closed
  $2\delta'$-neighbourhood of $F(L)$ lies in $\Omega'$ and that, whenever
  $\operatorname{dist}(x',F(L))\le\delta'$ and $\abs{x'-y'}\le\delta'$, the matrix
  \begin{equation*}
    T(x',y')=\int_0^1DG\bigl(y'+t(x'-y')\bigr)\,dt
  \end{equation*}
  is invertible with $\abs{\det T}\ge c>0$ and $\norm{T^{-1}}\le C$; this is
  possible since $T(x',x')=DG(x')$ and $F(L)$ is compact. Then
  $G(x')-G(y')=T(x',y')(x'-y')$. By uniform continuity of $F$ on a compact
  neighbourhood of $L$, choose $\delta>0$ such that $x,y\in L$,
  $\abs{x-y}\le2\delta$ imply $\abs{F(x)-F(y)}<\delta'$.

  Let $\mu\in\CIc(\R^n)$ equal $1$ on $\abs z\le\delta$ and vanish for
  $\abs z\ge2\delta$, and split $A=A_1+A_2$ with
  $K_{A_1}=\mu(x-y)K_A$, $K_{A_2}=(1-\mu(x-y))K_A$. By
  \cref{microlocal:lem:off-diagonal}, $A_2\in\Psido{-\infty}(\R^n)$, so
  $K_{A_2}$ is smooth, and it has compact support in $\Omega\times\Omega$.
  Hence $K_{(A_2)_F}$ is smooth with compact support in $\Omega'\times\Omega'$,
  in particular it satisfies the estimates characterizing
  $\Psido{-\infty}(\R^n)$, and $(A_2)_F\in\Psido{-\infty}(\R^n)$.

  The kernel of $A_1$ is $I(b)$ with
  $b(x,y,\xi)=\mu(x-y)\psi(x)\psi(y)a(x,\xi)\in S^m(\R^{2n};\R^n)$, which
  vanishes unless $x,y\in L$ and $\abs{x-y}\le2\delta$. If $b\in
  S^{-\infty}$ the integral converges absolutely and the substitution
  $\xi=S(x',y')\eta$, $S=(T^t)^{-1}$, gives
  \begin{equation}\label{microlocal:eq:AF-kernel}
  \begin{gathered}
    K_{(A_1)_F}(x',y')=\int e^{i(x'-y')\cdot\eta}c(x',y',\eta)\,\dbar\eta,\\
    c(x',y',\eta)=b\bigl(G(x'),G(y'),S(x',y')\eta\bigr)
    \frac{\abs{\det DG(y')}}{\abs{\det T(x',y')}},
  \end{gathered}
  \end{equation}
  where $c$ is extended by zero outside the compact set
  $\{(x',y');\ G(x'),G(y')\in L,\ \abs{G(x')-G(y')}\le2\delta\}$; on this set
  $\abs{x'-y'}<\delta'$, so $c$ is smooth. For general $b$ we apply this to
  $\chi(\eps\xi)b$ and let $\eps\to0$: both sides of
  \eqref{microlocal:eq:AF-kernel} converge in $\cSp(\R^{2n})$, by the
  continuity of $I$ on amplitudes and of the pull-back on compactly supported
  distributions. Thus $(A_1)_F=\Op(c)$, provided we show $c\in
  S^m(\R^{2n};\R^n)$.

  Write $f(x',y')$ for the smooth factor in $c$. By the chain rule,
  $D^\alpha_{x'}D^\beta_{y'}D^\gamma_\eta\bigl[b(G(x'),G(y'),S\eta)\bigr]$ is
  a finite sum of terms
  \begin{equation*}
    M(x',y')\,\eta^\nu\,
    \bigl(D^{\alpha'}_xD^{\beta'}_yD^\mu_\xi b\bigr)\bigl(G(x'),G(y'),S\eta\bigr),
    \qquad\abs\nu\le\abs\mu-\abs\gamma,
  \end{equation*}
  with smooth coefficients $M$: each $\eta$-derivative produces one
  $\xi$-derivative of $b$ and no power of $\eta$, while an $x'$- or
  $y'$-derivative falling on $S(x',y')\eta$ produces one $\xi$-derivative and
  at most one factor of $\eta$. On the support, $\ang{S\eta}$ is comparable
  to $\ang\eta$, so each term is bounded by
  $C\ang\eta^{\abs\nu}\ang\eta^{m-\abs\mu}\le C\ang\eta^{m-\abs\gamma}$. With
  Leibniz' formula for the factor $f$, this shows $c\in S^m(\R^{2n};\R^n)$.

  Thus $A_F=(A_1)_F+(A_2)_F\in\Psido m(\R^n)$. By
  \cref{ch:euclidean}, $\sigma_m(A_F)$ is represented by
  $c(x',x',\eta)$. Since $T(x',x')=DG(x')$,
  \begin{equation*}
    c(x',x',\eta)=\mu(0)\psi(G(x'))^2a\bigl(G(x'),(DG(x')^t)^{-1}\eta\bigr)
    =a\bigl(x,DF(x)^t\eta\bigr),\quad x=G(x'),
  \end{equation*}
  using $DG(x')^{-1}=DF(x)$ and $\psi a=a$. With $(x',\eta)=\chi_F(x,\xi)$ this
  is $a(x,\xi)$, proving the statement about principal symbols. The same
  chain rule estimates show that composition with $\chi_F^{-1}$ maps symbols
  of order $m-1$ supported over $L$ to symbols of order $m-1$, so the class
  is well defined.

  For the essential support we use \cref{microlocal:lem:ess-amplitude}. Let
  $(\bar x',\bar\eta)\in\Omega'\times(\R^n\setminus0)$ with
  $\chi_F^{-1}(\bar x',\bar\eta)=(\bar x,\bar\xi)\notin\WFp(A)$, so $a$, and
  hence $b(x,y,\xi)$ for $y$ arbitrary, is rapidly decreasing for
  $(x,\xi)\in\Gamma_\eps(\bar x,\bar\xi)$. Since $S(\bar x',\bar x')\bar\eta=\bar\xi$
  and $S$, $G$ are continuous, for $(x',y')$ close to $(\bar x',\bar x')$ and
  $\eta\in V_{\eps'}(\bar\eta)$ with $\eps'$ small we have
  $(G(x'),S(x',y')\eta)\in\Gamma_\eps(\bar x,\bar\xi)$; as
  $\ang{S\eta}\ge c\ang\eta$, $c$ is rapidly decreasing there. Hence $(\bar
  x',\bar x',\bar\eta)\notin\conesupp(c)$ and $(\bar x',\bar\eta)\notin\WFp(A_F)$.
  Points over $\R^n\setminus F(K)$ are not in $\WFp(A_F)$, since
  $\sigma_L(A_F)$ vanishes there. Thus $\WFp(A_F)\subset\chi_F(\WFp(A))$.
  Since $A=(A_F)_G$, the reverse inclusion follows by symmetry. The
  statement about characteristic varieties follows from that about principal
  symbols, since $\chi_F$ is homogeneous of degree one in the fibres and
  satisfies $c^{-1}\abs\xi\le\abs{(DF(x)^t)^{-1}\xi}\le c\abs\xi$ locally
  uniformly in $x$; so it maps the sets $\Gamma_\eps(z)\cap\{\abs\xi\ge
  1/\eps\}$ into and onto sets of the same type.
\end{proof}

This invariance shows that the principal symbol, essential support and
characteristic variety of a pseudodifferential operator are invariantly
defined on the cotangent bundle; this is the basis of the extension of the
calculus to manifolds in \cref{ch:manifolds}. It also gives us the freedom to
choose local coordinates adapted to a particular problem.

% ======================================================================
\section{The wavefront set}\label{microlocal:sec:wf}

We first define the wavefront set for distributions of compact support and
then localize.

\begin{definition}\label{microlocal:def:wf-compact}
  For $u\in\CmIc(\R^n)$,
  \begin{equation*}
    \WF(u)=\bigcap\bigl\{\Char_0(A);\ A\in\Psido0(\R^n),\ Au\in\CI(\R^n)\bigr\}.
  \end{equation*}
\end{definition}

Thus $z\notin\WF(u)$ if and only if there is $A\in\Psido0(\R^n)$, elliptic at
$z$, with $Au\in\CI(\R^n)$. As an intersection of closed conic sets,
$\WF(u)\subset\Tz{\R^n}$ is closed and conic. The wavefront set refines the
singular support. Let $\pi:\Tz{\R^n}\to\R^n$ be the projection.

\begin{proposition}\label{microlocal:prop:wf-singsupp}
  For $u\in\CmIc(\R^n)$, $\pi(\WF(u))=\singsupp(u)$.
\end{proposition}
\begin{proof}
  If $\bar x\notin\singsupp u$ there is $\phi\in\CIc(\R^n)$ with
  $\phi(\bar x)\ne0$ and $\phi u\in\CI$. As a multiplication operator
  $\phi\in\Psido0(\R^n)$ is elliptic at $(\bar x,\bar\xi)$ for every
  $\bar\xi\ne0$, so $(\bar x,\bar\xi)\notin\WF(u)$. Thus
  $\pi(\WF(u))\subset\singsupp(u)$.

  Conversely suppose $(\bar x,\omega)\notin\WF(u)$ for every
  $\omega\in\Sphere^{n-1}$. For each $\omega$ there is $A_\omega\in\Psido0$
  with $A_\omega u\in\CI$ and $\abs{\sigma_L(A_\omega)}\ge\eps_\omega$ on
  $\Gamma_{\eps_\omega}(\bar x,\omega)\cap\{\abs\xi\ge1/\eps_\omega\}$. By
  compactness of the sphere there are finitely many $\omega_j$ such that the
  $V_{\eps_j}(\omega_j)$, $\eps_j=\eps_{\omega_j}$, cover $\R^n\setminus0$. Let
  $A_j=A_{\omega_j}$, $\eps_0=\min_j\eps_j$, $R_0=\max_j1/\eps_j$ and
  \begin{equation*}
    B=\sum_jA_j^*A_j\in\Psido0(\R^n).
  \end{equation*}
  Since $A_ju\in\CI\cap\cSp$ and $A_j^*$ is pseudolocal
  (\cref{microlocal:prop:pseudolocal}), $Bu\in\CI(\R^n)$. Moreover
  $\sigma_L(B)=\sum_j\abs{\sigma_L(A_j)}^2+r$ with $r\in S^{-1}$, and
  $\sum_j\abs{\sigma_L(A_j)}^2\ge\eps_0^2$ when $\abs{x-\bar x}<\eps_0$ and
  $\abs\xi\ge R_0$. Let $\phi\in\CIc(\R^n)$, $0\le\phi\le1$, with support in
  $\{\abs{x-\bar x}<\eps_0\}$ and $\phi=1$ near $\bar x$. Then
  $\sigma_L(B+(1-\phi))=\sum_j\abs{\sigma_L(A_j)}^2+(1-\phi)+r$ and, for
  $\abs\xi\ge R_0$, the sum of the first two terms is at least
  $\min(\eps_0^2,1)$ (consider separately $\abs{x-\bar x}<\eps_0$ and
  $\abs{x-\bar x}\ge\eps_0$). As $\abs r\le C\ang\xi^{-1}$, the operator
  $B+(1-\phi)$ is globally elliptic, so has a parametrix
  $Q\in\Psido0(\R^n)$:
  \begin{equation*}
    \Id=QB+Q(1-\phi)-E,\qquad E\in\Psido{-\infty}(\R^n).
  \end{equation*}
  Let $\psi\in\CIc(\R^n)$ with $\psi(\bar x)\ne0$ and $\supp\psi$ contained in
  the interior of $\{\phi=1\}$. Then
  $\psi u=\psi QBu+\psi Q(1-\phi)u-\psi Eu$. The first term is smooth by
  pseudolocality, the second because $\psi Q(1-\phi)\in\Psido{-\infty}$ by
  \cref{microlocal:prop:kernel-singsupp}, and the third because
  $E\in\Psido{-\infty}$. Hence $\bar x\notin\singsupp u$.
\end{proof}

\begin{lemma}\label{microlocal:lem:wf-cone}
  Let $u\in\CmIc(\R^n)$ and $z\notin\WF(u)$. Then there is an open conic
  neighbourhood $\Gamma$ of $z$ such that $Cu\in\CI(\R^n)$ for every
  $C\in\Psido{m}(\R^n)$, $m\in\R$, with $\WFp(C)\subset\Gamma$.
\end{lemma}
\begin{proof}
  Let $A\in\Psido0$ be elliptic at $z$ with $Au\in\CI$ and let $Q$ be a
  microlocal parametrix for $A$ at $z$
  (\cref{microlocal:prop:parametrix}), so $QA=\Id+R$ with $z\notin\WFp(R)$.
  Let $\Gamma$ be the complement of $\WFp(R)$, an open conic neighbourhood of
  $z$, and let $\psi\in\CIc(\R^n)$ equal $1$ near $\supp u$. If
  $\WFp(C)\subset\Gamma$ then
  \begin{equation*}
    Cu=CQ(Au)-CR\psi u.
  \end{equation*}
  The first term is smooth by pseudolocality. Since
  $\WFp(CR)\subset\WFp(C)\cap\WFp(R)=\emptyset$,
  $CR\psi\in\Psido{-\infty}$ by \cref{microlocal:lem:localize}, so the second
  term is smooth too.
\end{proof}

In particular, $\WF(u+v)\subset\WF(u)\cup\WF(v)$ for $u,v\in\CmIc(\R^n)$:
if $z$ is in neither set, take $C=G_{z,\eps}$ with $\overline{\Gamma_\eps(z)}$
inside both of the cones given by the lemma.

\begin{proposition}\label{microlocal:prop:microlocal-compact}
  If $u\in\CmIc(\R^n)$, $A\in\Psido m(\R^n)$ and $\phi\in\CIc(\R^n)$ then
  \begin{equation*}
    \WF(\phi Au)\subset\WFp(A)\cap\WF(u).
  \end{equation*}
\end{proposition}
\begin{proof}
  Let $\psi\in\CIc(\R^n)$ equal $1$ near $\supp u$.

  Suppose first $z\notin\WFp(A)$. Since $\WFp(A)$ is closed and conic, there
  is $\eps>0$ with $\overline{\Gamma_\eps(z)}\cap\WFp(A)=\emptyset$; let
  $C=G_{z,\eps}$, which is elliptic at $z$. Then
  $\WFp(C\phi A)\subset\WFp(C)\cap\WFp(A)=\emptyset$, and $C\phi A=\psi'C\phi
  A$ for a suitable $\psi'\in\CIc$ by \eqref{microlocal:eq:G-cutoff}, so
  $C\phi A\in\Psido{-\infty}$ by \cref{microlocal:lem:localize}. Hence
  $C(\phi Au)\in\CI$ and $z\notin\WF(\phi Au)$.

  Suppose next $z\notin\WF(u)$ and let $\Gamma$ be as in
  \cref{microlocal:lem:wf-cone}. Choose $\eps$ with
  $\overline{\Gamma_\eps(z)}\subset\Gamma$ and let $D=G_{z,\eps}$, so
  $\WFp(D)\subset\Gamma$ and $Du\in\CI$; and let $C=G_{z,\eps/4}$, which is
  elliptic at $z$ with $\WFp(C)\subset\overline{\Gamma_{\eps/4}(z)}$, a set
  disjoint from $\WFp(\Id-D)$ by \eqref{microlocal:eq:G-cutoff}. Then
  \begin{equation*}
    C\phi Au=C\phi A(Du)+C\phi A(\Id-D)\psi u.
  \end{equation*}
  The first term is smooth by pseudolocality, as $Du\in\CI\cap\cSp$. In
  the second term,
  $\WFp(C\phi A(\Id-D))\subset\WFp(C)\cap\WFp(\Id-D)=\emptyset$, so
  $C\phi A(\Id-D)\psi\in\Psido{-\infty}$ by \cref{microlocal:lem:localize}.
  Hence $C(\phi Au)\in\CI$ and $z\notin\WF(\phi Au)$.
\end{proof}

Taking $A=\Id$ we see that $\WF(\phi u)\subset\WF(u)$ for $\phi\in\CIc$ and
$u\in\CmIc$, and $\WF(\phi u)=\WF(u)$ if $\phi=1$ near $\supp u$. This makes
the following extension consistent.

\begin{definition}\label{microlocal:def:wf}
  For $u\in\cSp(\R^n)$,
  \begin{equation*}
    \WF(u)=\bigcup_{\phi\in\CIc(\R^n)}\WF(\phi u)\subset\Tz{\R^n}.
  \end{equation*}
\end{definition}

\begin{proposition}\label{microlocal:prop:wf-local}
  Let $u,v\in\cSp(\R^n)$.
  \begin{enumerate}[label=\textup{(\roman*)}]
  \item If $\phi\in\CIc(\R^n)$ and $\phi(\bar x)\ne0$ then
    $\WF(u)\cap\pi^{-1}(\bar x)=\WF(\phi u)\cap\pi^{-1}(\bar x)$.
  \item $\WF(u)$ is closed and conic, and $\pi(\WF(u))=\singsupp(u)$.
  \item $\WF(u+v)\subset\WF(u)\cup\WF(v)$.
  \end{enumerate}
\end{proposition}
\begin{proof}
  (i) One inclusion is the definition. Conversely let $z=(\bar x,\bar\xi)\notin
  \WF(\phi u)$ and $\psi\in\CIc(\R^n)$. Choose $\phi'\in\CIc$ supported in
  $\{\phi\ne0\}$ with $\phi'=1$ near $\bar x$, and write
  $\psi u=\theta(\phi u)+\psi(1-\phi')u$ with
  $\theta=\psi\phi'/\phi\in\CIc$. Then $z\notin\WF(\theta\phi u)$ by
  \cref{microlocal:prop:microlocal-compact}, while $\psi(1-\phi')u$ vanishes
  near $\bar x$, so its wavefront set does not meet $\pi^{-1}(\bar x)$ by
  \cref{microlocal:prop:wf-singsupp}. Hence $z\notin\WF(\psi u)$.

  (ii) By (i), $\WF(u)\cap\pi^{-1}(U)=\WF(\phi u)\cap\pi^{-1}(U)$ on the open
  set $U=\{\phi\ne0\}$, so $\WF(u)$ is closed, being locally closed; it is
  clearly conic. Also $\bar x\notin\singsupp u$ if and only if
  $\bar x\notin\singsupp(\phi u)$ for some (any) $\phi$ as in (i), so the
  second statement follows from \cref{microlocal:prop:wf-singsupp}.

  (iii) follows from (i) and the corresponding statement for $\CmIc(\R^n)$.
\end{proof}

\begin{theorem}[Microlocality]\label{microlocal:thm:microlocality}
  If $A\in\Psido m(\R^n)$ and $u\in\cSp(\R^n)$ then
  \begin{equation*}
    \WF(Au)\subset\WFp(A)\cap\WF(u).
  \end{equation*}
\end{theorem}
\begin{proof}
  Let $\phi\in\CIc(\R^n)$ and choose $\psi\in\CIc(\R^n)$ equal to $1$ near
  $\supp\phi$. Then $\phi Au=\phi A(\psi u)+\phi A(1-\psi)u$, and the
  second term is smooth by \cref{microlocal:prop:kernel-singsupp}. By
  \cref{microlocal:prop:microlocal-compact},
  $\WF(\phi Au)\subset\WFp(A)\cap\WF(\psi u)\subset\WFp(A)\cap\WF(u)$. Taking
  the union over $\phi$ gives the result.
\end{proof}

Combined with microlocal parametrices this gives the basic regularity
result.

\begin{theorem}[Microlocal elliptic regularity]\label{microlocal:thm:elliptic-regularity}
  If $A\in\Psido m(\R^n)$ and $u\in\cSp(\R^n)$ then
  \begin{equation*}
    \WF(u)\subset\WF(Au)\cup\Char_m(A).
  \end{equation*}
\end{theorem}
\begin{proof}
  Let $z\notin\Char_m(A)$ and let $B$ be a microlocal parametrix at $z$, so
  $BA=\Id+R$ with $z\notin\WFp(R)$. Then $u=B(Au)-Ru$, and by
  \cref{microlocal:thm:microlocality} and
  \cref{microlocal:prop:wf-local}(iii),
  $\WF(u)\subset\WF(Au)\cup\WFp(R)$, which does not contain $z$ if
  $z\notin\WF(Au)$.
\end{proof}

In particular, if $A$ is elliptic at every point (that is,
$\Char_m(A)=\emptyset$), then $\WF(u)=\WF(Au)$. For instance $\WF(u)$
is unchanged by the Fourier multipliers $\ang D^s$.

% ======================================================================
\section{The wavefront set and the Fourier transform}\label{microlocal:sec:fourier}

We now give the more classical characterization of the wavefront set, in
terms of the Fourier transform. We use the following standard facts about
$u\in\CmIc(\R^n)$, which follow from \eqref{microlocal:eq:finite-order} and
the Fourier inversion formula by approximating the integrals below by
Riemann sums in $\cS(\R^n)$: $\hat u(\xi)=u(e^{-ix\cdot\xi})$ is a smooth
function with $\abs{\hat u(\xi)}\le C\ang\xi^k$; for $\phi\in\cS(\R^n)$,
\begin{equation}\label{microlocal:eq:parseval}
  u(\phi)=\int\hat u(\xi)\,\hat\phi(-\xi)\,\dbar\xi,\qquad
  \widehat{\phi u}(\xi)=\int\hat\phi(\xi-\eta)\,\hat u(\eta)\,\dbar\eta .
\end{equation}
Conversely, a continuous polynomially bounded function is the Fourier
transform of a unique element of $\cSp(\R^n)$.

For $v\in\CmIc(\R^n)$ let
\begin{equation*}
  \Sigma(v)=\{\xi\in\R^n\setminus0;\ (x,\xi)\in\WF(v)\text{ for some }x\}.
\end{equation*}
Since $\WF(v)$ is closed, conic and lies over the compact set $\supp v$,
$\WF(v)\cap(\R^n\times\Sphere^{n-1})$ is compact, so $\Sigma(v)$ is a closed
cone in $\R^n\setminus0$, and $\WF(v)\subset\supp v\times\Sigma(v)$.

\begin{lemma}\label{microlocal:lem:fourier-global}
  Let $v\in\CmIc(\R^n)$.
  \begin{enumerate}[label=\textup{(\roman*)}]
  \item If $\Gamma\subset\R^n\setminus0$ is a closed cone with
    $\Gamma\cap\Sigma(v)=\emptyset$, then $\hat v$ is rapidly decreasing on
    $\Gamma$.
  \item If $\hat v$ is rapidly decreasing on an open cone $V$ then
    $\WF(v)\cap(\R^n\times V)=\emptyset$.
  \item If $\bar x\in\R^n$ and $V$ is an open cone containing
    $\{\xi;\ (\bar x,\xi)\in\WF(v)\}$, there is $r>0$ such that
    $\Sigma(\phi v)\subset V$ for all $\phi\in\CIc(\R^n)$ supported in
    $\{\abs{x-\bar x}<r\}$.
  \end{enumerate}
\end{lemma}
\begin{proof}
  (i) By compactness there is an open cone $\Gamma'\supset\Gamma$ with
  $\overline{\Gamma'}\cap\Sigma(v)=\emptyset$ (the distance between
  $\Gamma\cap\Sphere^{n-1}$ and the compact set $\Sigma(v)\cap\Sphere^{n-1}$ is
  positive). Choose $g\in S^0(\R^n)$, a function of $\xi$ only, with $g=1$ on
  $\Gamma\cap\{\abs\xi\ge1\}$ and $\supp g\subset\Gamma'$; for instance $g=
  h(\xi/\abs\xi)(1-\chi(\xi))$ with $h\in\CI(\Sphere^{n-1})$. Then
  $g(D)=\Op(g)$ has $\WFp(g(D))\subset\R^n\times\overline{\Gamma'}$, which is
  disjoint from $\WF(v)$, so $g(D)v\in\CI$ by
  \cref{microlocal:thm:microlocality}. Let $\psi_1,\psi_2\in\CIc(\R^n)$ with
  $\psi_1=1$ near $\supp v$ and $\psi_2=1$ near $\supp\psi_1$. Then
  \begin{equation*}
    g(D)v=\psi_2g(D)v+(1-\psi_2)g(D)\psi_1v,
  \end{equation*}
  where the first term is in $\CIc$ and the second is in $\cS$ by
  \cref{microlocal:prop:kernel-singsupp} and
  \cref{microlocal:lem:residual-compact}. Thus $g\hat v=\widehat{g(D)v}\in\cS$,
  and $\hat v$ is rapidly decreasing on $\Gamma$.

  (ii) Let $z=(\bar x,\bar\xi)$ with $\bar\xi\in V$, and choose $\eps>0$ with
  $\overline{V_\eps(\bar\xi)}\subset V$ and $g\in S^0(\R^n)$ a function of
  $\xi$ with $g=1$ on $V_{\eps/2}(\bar\xi)\cap\{\abs\xi\ge1\}$, supported in
  $V_\eps(\bar\xi)$. Then $g(D)$ is elliptic at $z$, and
  $\widehat{g(D)v}=g\hat v$ is rapidly decreasing, hence integrable together
  with $\xi^\alpha g\hat v$ for all $\alpha$, so $g(D)v\in\CI$ by Fourier
  inversion. Thus $z\notin\WF(v)$.

  (iii) Otherwise there are $x_j\to\bar x$ and unit vectors $\xi_j\notin V$
  with $(x_j,\xi_j)\in\WF(v)$; a subsequence has $\xi_j\to\xi$, and
  $(\bar x,\xi)\in\WF(v)$ by closedness, $\xi\notin V$ since the complement of
  $V$ is closed, a contradiction. Since $\WF(\phi v)\subset\WF(v)\cap
  (\supp\phi\times\R^n)$, the claim follows.
\end{proof}

\begin{proposition}[Fourier characterization]\label{microlocal:prop:fourier-char}
  Let $u\in\cSp(\R^n)$ and $(\bar x,\bar\xi)\in\Tz{\R^n}$. Then
  $(\bar x,\bar\xi)\notin\WF(u)$ if and only if there exist
  $\phi\in\CIc(\R^n)$ with $\phi(\bar x)\ne0$ and $\eps>0$ such that
  $\widehat{\phi u}$ is rapidly decreasing on $V_\eps(\bar\xi)$:
  \begin{equation*}
    \abs{\widehat{\phi u}(\xi)}\le C_N\ang\xi^{-N}\quad\text{for }
    \xi\in V_\eps(\bar\xi),\ N\in\N.
  \end{equation*}
  More precisely, if $K\subset\R^n$ is compact, $\Gamma\subset\R^n\setminus0$
  is a closed cone and $\WF(u)\cap(K\times\Gamma)=\emptyset$, then
  $\widehat{\phi u}$ is rapidly decreasing on $\Gamma$ for every
  $\phi\in\CIc(\R^n)$ with $\supp\phi\subset K$.
\end{proposition}
\begin{proof}
  If $\widehat{\phi u}$ is rapidly decreasing on $V_\eps(\bar\xi)$ then
  $(\bar x,\bar\xi)\notin\WF(\phi u)$ by
  \cref{microlocal:lem:fourier-global}(ii), and so $(\bar
  x,\bar\xi)\notin\WF(u)$ by \cref{microlocal:prop:wf-local}(i).

  For the last statement, $\WF(\phi u)\subset\WF(u)\cap(K\times\R^n)$, so
  $\Sigma(\phi u)\cap\Gamma=\emptyset$ and
  \cref{microlocal:lem:fourier-global}(i) applies. Finally, if $(\bar
  x,\bar\xi)\notin\WF(u)$, closedness of $\WF(u)$ and a compactness argument
  as in \cref{microlocal:lem:fourier-global}(iii) give $r,\eps>0$ with
  $\WF(u)\cap(K\times\overline{V_\eps(\bar\xi)})=\emptyset$ for
  $K=\{\abs{x-\bar x}\le r\}$, and any $\phi$ supported in $K$ with
  $\phi(\bar x)\ne0$ will do.
\end{proof}

\begin{example}\label{microlocal:ex:delta}
  Since $\hat\delta=1$ and $\singsupp\delta=\{0\}$,
  $\WF(\delta)=\{0\}\times(\R^n\setminus0)$: indeed
  $\widehat{\phi\delta}=\phi(0)$ is not rapidly decreasing in any direction if
  $\phi(0)\ne0$. More generally, if $u\in\CmIc(\R^n)$ is not smooth then
  $\Sigma(u)\ne\emptyset$, since otherwise $\hat u\in\cS$ by
  \cref{microlocal:lem:fourier-global}(i).
\end{example}

\begin{corollary}\label{microlocal:cor:wf-conj}
  For $u\in\cSp(\R^n)$, $\WF(\bar u)=\WF(u)^-$.
\end{corollary}
\begin{proof}
  $\widehat{\phi\bar u}(\xi)=\overline{\widehat{\bar\phi u}(-\xi)}$, so this
  follows from \cref{microlocal:prop:fourier-char} applied with $\bar\phi$.
\end{proof}

% ======================================================================
\section{Invariance under diffeomorphisms}\label{microlocal:sec:diffeo}

For a diffeomorphism $F:\Omega\to\Omega'$ between open subsets of $\R^n$ we
use the map $\chi_F(x,\xi)=(F(x),(DF(x)^t)^{-1}\xi)$ of
\cref{microlocal:sec:coordinates}. Its inverse is
\begin{equation*}
  \chi_F^{-1}(y,\eta)=\bigl(G(y),DF(G(y))^t\eta\bigr),\qquad G=F^{-1},
\end{equation*}
the cotangent map $\eta\mapsto DF(x)^t\eta$, $x=G(y)$, which we denote by
$F^*$ when applied to subsets: for $W\subset\Tz{\Omega'}$,
\begin{equation*}
  F^*W=\{(x,DF(x)^t\eta);\ (F(x),\eta)\in W\}.
\end{equation*}

\begin{theorem}\label{microlocal:thm:wf-diffeo}
  If $F:\Omega\to\Omega'$ is a diffeomorphism between open subsets of $\R^n$
  and $u\in\CmIc(\Omega')$, then
  \begin{equation*}
    \WF(F^*u)=F^*\WF(u)=\{(x,DF(x)^t\eta);\ (F(x),\eta)\in\WF(u)\}.
  \end{equation*}
\end{theorem}
\begin{proof}
  Let $(x,\xi)\in\Tz\Omega$ with $w=\chi_F(x,\xi)\notin\WF(u)$; we show
  $(x,\xi)\notin\WF(F^*u)$. Let $\Gamma$ be an open conic neighbourhood of $w$
  as in \cref{microlocal:lem:wf-cone}, choose $\eps>0$ with
  $\overline{\Gamma_\eps(w)}\subset\Gamma$ and $\{\abs{y-F(x)}\le\eps\}\subset
  \Omega'$, and let $\theta\in\CIc(\Omega')$ equal $1$ on
  $\{\abs{y-F(x)}\le\eps\}$. The operator
  $B=\theta G_{w,\eps}\theta\in\Psido0(\R^n)$ has kernel with compact support
  in $\Omega'\times\Omega'$, $\WFp(B)\subset\overline{\Gamma_\eps(w)}\subset\Gamma$
  and is elliptic at $w$. So $Bu\in\CIc(\Omega')$.

  By \cref{microlocal:thm:coordinates}, applied to the diffeomorphism
  $G:\Omega'\to\Omega$, the operator $\tilde B=B_G=F^*BG^*$ is in
  $\Psido0(\R^n)$ and is elliptic at $\chi_G(w)=(x,\xi)$. Both
  $\tilde B$ (as an element of $\Psido0(\R^n)$ acting on $\cSp$) and
  $F^*\circ B\circ G^*$ are weakly continuous on $\CmIc(\Omega)$ and agree on
  $\CIc(\Omega)$, so by density they agree on $\CmIc(\Omega)$. Hence
  $\tilde B(F^*u)=F^*BG^*F^*u=F^*(Bu)\in\CI$, and $(x,\xi)\notin\WF(F^*u)$.

  We have shown $\WF(F^*u)\subset\chi_F^{-1}(\WF(u))=F^*\WF(u)$. Applying
  this to $G$ and $F^*u$, and using $G^*F^*u=u$, gives the reverse
  inclusion.
\end{proof}

If $L$ is an invertible linear map of $\R^n$ and $u\in\cSp(\R^n)$, then by
localization (\cref{microlocal:prop:wf-local}) the theorem also applies to
$L^*u$: $\WF(L^*u)=\{(x,L^t\eta);\ (Lx,\eta)\in\WF(u)\}$.

% ======================================================================
\section{The wavefront set of a kernel}\label{microlocal:sec:kernel}

The kernels of pseudodifferential operators form an interesting class of
distributions. Their wavefront sets are determined by, and determine, the
essential supports of the operators.

\begin{proposition}\label{microlocal:prop:wf-kernel}
  If $A\in\Psido m(\R^n)$ then
  \begin{equation*}
    \WF(K_A)=\{(x,x,\xi,-\xi)\in\Tz{\R^{2n}};\ (x,\xi)\in\WFp(A)\}.
  \end{equation*}
\end{proposition}
\begin{proof}
  By \cref{microlocal:prop:kernel-singsupp} and
  \cref{microlocal:prop:wf-local}(ii), $\WF(K_A)$ lies over the diagonal.
  Fix $\bar x\in\R^n$, let $\mu\in\CIc(\R^n)$ equal $1$ near $0$ and let
  $A_\mu$ have kernel $\mu(x-y)K_A$; by \cref{microlocal:lem:off-diagonal},
  $A-A_\mu\in\Psido{-\infty}$, so $a_\mu=\sigma_L(A_\mu)$ differs from
  $a=\sigma_L(A)$ by an element of $S^{-\infty}$. For $\phi\in\CIc(\R^n)$
  with $\phi(\bar x)\ne0$ consider
  \begin{equation*}
    k(x,y)=\phi(x)\mu(x-y)K_A(x,y)\in\CmIc(\R^{2n}),\qquad
    w=L^*k,\quad L(x,z)=(x,x-z),
  \end{equation*}
  so $w(x,z)=\phi(x)\mu(z)K_A(x,x-z)$, and $k=K_A$ near $(\bar x,\bar x)$.
  Since $K_{A_\mu}=\int e^{i(x-y)\cdot\xi}a_\mu(x,\xi)\,\dbar\xi$, the
  distribution $L^*K_{A_\mu}$ is the inverse Fourier transform of
  $a_\mu(x,\xi)$ in $\xi$, and so
  \begin{equation*}
    \hat w(\zeta,\xi)=\int e^{-ix\cdot\zeta}\phi(x)a_\mu(x,\xi)\,dx,
  \end{equation*}
  the Fourier transform being taken in $(x,z)\mapsto(\zeta,\xi)$. Since
  $\phi a_\mu$ is a symbol with compact support in $x$, integration by parts
  in $x$ gives
  \begin{equation}\label{microlocal:eq:kernel-ft}
    \abs{\hat w(\zeta,\xi)}\le C_N\ang\zeta^{-N}\ang\xi^m\quad\text{for all }N.
  \end{equation}
  Points of the wavefront set and the covectors transform under $L$ by
  \cref{microlocal:thm:wf-diffeo}: since $DL^t(\xi,\eta)=(\xi+\eta,-\eta)$,
  \begin{equation}\label{microlocal:eq:kernel-L}
    (x,0;\zeta,\eta')\in\WF(w)\iff
    (x,x;\zeta+\eta',-\eta')\in\WF(k).
  \end{equation}

  \emph{Inclusion $\subset$.} By \eqref{microlocal:eq:kernel-ft}, $\hat w$ is
  rapidly decreasing on every closed cone $\{\abs\xi\le c\abs\zeta\}$, so
  $\Sigma(w)\subset\{\zeta=0\}$ by \cref{microlocal:lem:fourier-global}(ii).
  Since $w$ is smooth for $z\ne0$, $\WF(w)\subset\{(x,0;0,\xi)\}$, and
  \eqref{microlocal:eq:kernel-L} gives $\WF(k)\subset\{(x,x;\xi,-\xi)\}$.

  Now suppose $(\bar x,\bar\xi)\notin\WFp(A)$. Then $a_\mu$ and all its
  derivatives are rapidly decreasing on $\Gamma_\eps(\bar x,\bar\xi)$ for some
  $\eps$ (\cref{microlocal:lem:conesupp-derivatives}); take
  $\supp\phi\subset\{\abs{x-\bar x}<\eps\}$. Then integration by parts gives
  $\abs{\hat w(\zeta,\xi)}\le C_N\ang\zeta^{-N}\ang\xi^{-N}$ for
  $\xi\in V_\eps(\bar\xi)$, which with \eqref{microlocal:eq:kernel-ft} shows
  that $\hat w$ is rapidly decreasing on the cone
  $\{(\zeta,\xi);\ \xi\in V_\eps(\bar\xi),\ \abs\zeta<\abs\xi\}$, a conic
  neighbourhood of $(0,\bar\xi)$. Hence $(\bar x,0;0,\bar\xi)\notin\WF(w)$,
  and by \eqref{microlocal:eq:kernel-L},
  $(\bar x,\bar x;\bar\xi,-\bar\xi)\notin\WF(k)$, so not in $\WF(K_A)$.

  \emph{Inclusion $\supset$.} Suppose
  $(\bar x,\bar x,\bar\xi,-\bar\xi)\notin\WF(K_A)$. By closedness there are
  $r,\theta>0$ such that $(x,x,\xi,-\xi)\notin\WF(K_A)$ for $\abs{x-\bar
  x}<r$, $\xi\in V_\theta(\bar\xi)$. Take $\supp\phi\subset\{\abs{x-\bar
  x}<r\}$ and $\phi=1$ near $\bar x$. Then, by
  \eqref{microlocal:eq:kernel-L} and the first part,
  $\Sigma(w)\subset\{(0,\xi);\ \xi\notin V_\theta(\bar\xi)\}$, which is
  disjoint from the closed cone
  $C=\{(\zeta,\xi)\ne0;\ \abs\zeta\le\abs\xi,\ \xi\in
  \overline{V_{\theta/2}(\bar\xi)}\}$. By
  \cref{microlocal:lem:fourier-global}(i), $\hat w$ is rapidly decreasing on
  $C$. By Fourier inversion in $\zeta$, for $\xi\in V_{\theta/2}(\bar\xi)$,
  \begin{equation*}
    \abs{\phi(x)a_\mu(x,\xi)}\le\int_{\abs\zeta\le\abs\xi}\abs{\hat w(\zeta,\xi)}\,\dbar\zeta
    +\int_{\abs\zeta>\abs\xi}\abs{\hat w(\zeta,\xi)}\,\dbar\zeta
    \le C_N\ang\xi^{n-N}+C_N\ang\xi^{m+n-N},
  \end{equation*}
  using \eqref{microlocal:eq:kernel-ft} in the second integral. So $a_\mu$,
  and hence $a$, is rapidly decreasing on $\Gamma_{r'}(\bar x,\bar\xi)$ for
  small $r'$, that is $(\bar x,\bar\xi)\notin\WFp(A)$.
\end{proof}

In particular $\WFp(A)$ and $\WF(K_A)$ determine each other, and
$\WFp(A)=\emptyset$ if and only if $K_A\in\CI(\R^{2n})$.

% ======================================================================
\section{The calculus of wavefront sets: pairing}\label{microlocal:sec:pairing}

We now wish to understand, in terms of wavefront sets, three fundamental
operations on distributions: pull-back $F^*u$ by a smooth map, push-forward
$F_*u$, and multiplication $u_1u_2$. None of them is defined on all
distributions. To analyse them we first discuss some more elementary
operations:
\begin{align*}
  &\text{pairing:}&&(u,v)\longmapsto\int u(x)v(x)\,dx,\\
  &\text{exterior product:}&&(u,v)\longmapsto(u\boxtimes v)(x,y)=u(x)v(y),\\
  &\text{restriction:}&&u(x,y)\longmapsto u(x,0),\\
  &\text{projection:}&&u(x,y)\longmapsto\int u(x,y)\,dy,\\
  &\text{invariance:}&&u\longmapsto F^*u,\ F\text{ a diffeomorphism}.
\end{align*}
The last was treated in \cref{microlocal:sec:diffeo}. Restriction and
diffeomorphisms are special cases of pull-back and projection is a special
case of push-forward; conversely, we shall express the three fundamental
operations in terms of the elementary ones. Throughout, all distributions
have compact support; by the local nature of the wavefront set this is not
a serious restriction.

The pairing $\int u_1u_2\,dx$ of two distributions is the first example of an
operation which requires a condition on wavefront sets. We give Melrose's
operator-theoretic definition; in \cref{microlocal:sec:products} it is
recovered as the integral of the product (\cref{microlocal:ex:pairing-product}).
We need a transposition identity.

\begin{lemma}\label{microlocal:lem:transpose}
  Let $D\in\Psido m(\R^n)$ and $u,v\in\CmIc(\R^n)$ with
  $\WFp(D)\cap\WF(u)=\emptyset$ and $\WFp(D^t)\cap\WF(v)=\emptyset$. Then
  $Du$ and $D^tv$ are smooth and, for any $\psi\in\CIc(\R^n)$ equal to $1$
  near $\supp u\cup\supp v$,
  \begin{equation*}
    v(\psi Du)=u(\psi D^tv).
  \end{equation*}
\end{lemma}
\begin{proof}
  Smoothness follows from \cref{microlocal:thm:microlocality}. Both sides
  are independent of the choice of $\psi$. Let $\psi_1\in\CIc(\R^n)$ equal $1$
  near $\supp u\cup\supp v$ and let $\psi=\psi_2$ equal $1$ near
  $\supp\psi_1$. The closed conic sets
  $W_1=\WFp(D)\cap(\supp\psi_1\times\R^n)$ and $W_2=\WF(u)$ are disjoint and
  $\pi(W_1)$ is bounded, so \cref{microlocal:lem:cutoff} gives
  $B\in\Psido0(\R^n)$ with $\WFp(B)\cap\WF(u)=\emptyset$ and
  $\WFp(\Id-B)\cap W_1=\emptyset$. Write $u=u_s+u_r$ with
  $u_s=\psi_1Bu$ and $u_r=\psi_1(\Id-B)u$. By microlocality $Bu$ is smooth, so
  $u_s\in\CIc(\R^n)$, and since $D^tv$ is smooth, the definition of $D$ on
  $\cSp$ gives
  \begin{equation*}
    v(\psi Du_s)=v(Du_s)=(D^tv)(u_s)=\int(D^tv)u_s\,dx=u_s(\psi D^tv).
  \end{equation*}
  For $u_r$, let $E=\psi_2D\psi_1(\Id-B)\psi_1$. Since
  $\WFp(D\psi_1)\subset W_1$, $\WFp(E)\subset W_1\cap\WFp(\Id-B)=\emptyset$, and
  by \cref{microlocal:lem:localize} $E\in\Psido{-\infty}(\R^n)$ with kernel
  $K_E\in\CIc(\R^{2n})$. As $u=\psi_1u$, $\psi_2Du_r=Eu$, while
  \begin{equation*}
    u_r(\psi_2D^tv)=u\bigl((\Id-B^t)\psi_1\psi_2D^tv\bigr)
    =u\bigl(\psi_1(\Id-B^t)\psi_1D^t\psi_2v\bigr)=u(E^tv),
  \end{equation*}
  using $u=\psi_1u$, $v=\psi_2v$ and $\psi_1\psi_2=\psi_1$. Finally
  $v(Eu)=u(E^tv)$ for $E$ with kernel in $\CIc(\R^{2n})$: both sides equal
  $(v\boxtimes u)(K_E)$, since $Eu(x)=u(K_E(x,\cdot))$,
  $E^tv(y)=v(K_E(\cdot,y))$, and the identity
  $v_x(u_y(\Phi(x,y)))=u_y(v_x(\Phi(x,y)))$ holds for $\Phi=\phi\boxtimes\phi'$,
  hence by continuity and density for all $\Phi\in\cS(\R^{2n})$. Adding the
  two identities proves the lemma.
\end{proof}

\begin{proposition}[Pairing]\label{microlocal:prop:pairing}
  Let $u_1,u_2\in\CmIc(\R^n)$ satisfy
  \begin{equation}\label{microlocal:eq:pairing-cond}
    (x,\xi)\in\WF(u_1)\Longrightarrow(x,-\xi)\notin\WF(u_2),
    \quad\text{i.e.}\quad\WF(u_1)\cap\WF(u_2)^-=\emptyset.
  \end{equation}
  Then there exists $A\in\Psido0(\R^n)$ with
  \begin{equation}\label{microlocal:eq:pairing-A}
    \WFp(A)\cap\WF(u_1)=\emptyset,\qquad\WFp(\Id-A)\cap\WF(u_2)^-=\emptyset,
  \end{equation}
  and for any such $A$ and any $\psi\in\CIc(\R^n)$ equal to $1$ near $\supp
  u_1\cup\supp u_2$ the number
  \begin{equation*}
    P(u_1,u_2)=u_2(\psi Au_1)+u_1\bigl(\psi(\Id-A)^tu_2\bigr)
  \end{equation*}
  is independent of $A$ and $\psi$. The pairing $P$ is bilinear, symmetric,
  and $P(u_1,u_2)=u_2(u_1)=\int u_1u_2\,dx$ if $u_1\in\CIc(\R^n)$ (and similarly
  if $u_2\in\CIc$).
\end{proposition}
\begin{proof}
  Existence of $A$ follows from \cref{microlocal:lem:cutoff} with
  $W_1=\WF(u_2)^-$, $W_2=\WF(u_1)$. By microlocality and
  $\WFp((\Id-A)^t)=\WFp(\Id-A)^-$ (\cref{microlocal:prop:ess}), $Au_1$ and
  $(\Id-A)^tu_2$ are smooth, so $P(u_1,u_2)$ is defined; it is clearly
  independent of $\psi$. If $A'$ is another choice, $D=A-A'$ satisfies
  $\WFp(D)\cap\WF(u_1)=\emptyset$, and since $D=(\Id-A')-(\Id-A)$ also
  $\WFp(D)\cap\WF(u_2)^-=\emptyset$, that is
  $\WFp(D^t)\cap\WF(u_2)=\emptyset$. The difference of the two expressions is
  $u_2(\psi Du_1)-u_1(\psi D^tu_2)$, which vanishes by
  \cref{microlocal:lem:transpose}.

  For linearity in the first argument, suppose $(u_1,u_2)$ and $(u_1',u_2)$
  both satisfy \eqref{microlocal:eq:pairing-cond}. Then
  $(\WF(u_1)\cup\WF(u_1'))\cap\WF(u_2)^-=\emptyset$, and an $A$ obtained from
  \cref{microlocal:lem:cutoff} with $W_2=\WF(u_1)\cup\WF(u_1')$ satisfies
  \eqref{microlocal:eq:pairing-A} for the three pairs $(u_1,u_2)$,
  $(u_1',u_2)$ and $(u_1+u_1',u_2)$, since
  $\WF(u_1+u_1')\subset\WF(u_1)\cup\WF(u_1')$; with this $A$ the expressions are
  visibly additive. Linearity in the second argument is the same. For symmetry note that $A'=\Id-A^t$
  satisfies \eqref{microlocal:eq:pairing-A} for the pair $(u_2,u_1)$, and the
  corresponding expression is $u_1(\psi(\Id-A^t)u_2)+u_2(\psi Au_1)$. If
  $u_1\in\CIc$ we may take $A=\Id$, and if $u_2\in\CIc$, $A=0$.
\end{proof}

The \emph{sesquilinear} pairing $\pair{u_1}{u_2}=\int u_1\overline{u_2}\,dx$
is then defined as $P(u_1,\overline{u_2})$. Since
$\WF(\overline{u_2})=\WF(u_2)^-$ (\cref{microlocal:cor:wf-conj}), the
condition becomes simply $\WF(u_1)\cap\WF(u_2)=\emptyset$, and
$\pair{u_1}{u_2}=\pair{Au_1}{u_2}+\pair{u_1}{(\Id-A^*)u_2}$ for any $A$ with
$\WFp(A)\cap\WF(u_1)=\emptyset=\WFp(\Id-A)\cap\WF(u_2)$.

% ======================================================================
\section{Exterior products, restriction and projection}\label{microlocal:sec:elementary}

\subsection{Exterior product}
The exterior product $u_1\boxtimes u_2$ of $u_1\in\CmIc(\R^p)$ and
$u_2\in\CmIc(\R^k)$ is always defined (\cref{ch:distributions}); it lies in
$\CmIc(\R^{p+k})$, $(u_1\boxtimes u_2)(\Phi)=u_1(u_2(\Phi(x,\cdot)))$, and
its Fourier transform is
\begin{equation*}
  \widehat{u_1\boxtimes u_2}(\xi,\eta)=\hat u_1(\xi)\,\hat u_2(\eta).
\end{equation*}

\begin{proposition}\label{microlocal:prop:exterior}
  For $u_1\in\CmIc(\R^p)$ and $u_2\in\CmIc(\R^k)$,
  \begin{multline*}
    \WF(u_1\boxtimes u_2)\subset\bigl(\WF(u_1)\times\WF(u_2)\bigr)
    \cup\{(x,y,0,\eta);\ x\in\supp u_1,\ (y,\eta)\in\WF(u_2)\}\\
    \cup\{(x,y,\xi,0);\ (x,\xi)\in\WF(u_1),\ y\in\supp u_2\},
  \end{multline*}
  where $\WF(u_1)\times\WF(u_2)$ means
  $\{(x,y,\xi,\eta);(x,\xi)\in\WF(u_1),(y,\eta)\in\WF(u_2)\}$.
\end{proposition}
\begin{proof}
  Let $(\bar x,\bar y,\bar\xi,\bar\eta)$ be outside the right side. We localize
  with $\phi_1(x)\phi_2(y)$, $\phi_1(\bar x)\ne0\ne\phi_2(\bar y)$, using
  $\phi_1\phi_2(u_1\boxtimes u_2)=(\phi_1u_1)\boxtimes(\phi_2u_2)$, whose Fourier
  transform is $\widehat{\phi_1u_1}(\xi)\widehat{\phi_2u_2}(\eta)$. If
  $\bar x\notin\supp u_1$ or $\bar y\notin\supp u_2$ we can make this zero.
  Otherwise one of the following holds.

  (a) $\bar\xi\ne0$ and $(\bar x,\bar\xi)\notin\WF(u_1)$. Choose $\phi_1$ by
  \cref{microlocal:prop:fourier-char} with $\widehat{\phi_1u_1}$ rapidly
  decreasing on $V_\eps(\bar\xi)$. On the conic neighbourhood
  $\{(\xi,\eta);\ \xi\in V_\eps(\bar\xi),\ \abs\eta<C\abs\xi\}$ of
  $(\bar\xi,\bar\eta)$, where $C>\abs{\bar\eta}/\abs{\bar\xi}$, we have
  $\ang{(\xi,\eta)}\le(1+C)\ang\xi$ and
  $\abs{\widehat{\phi_2u_2}(\eta)}\le C'\ang\xi^M$, so the product is rapidly
  decreasing.

  (b) $\bar\eta\ne0$ and $(\bar y,\bar\eta)\notin\WF(u_2)$: the same with the
  factors exchanged.

  These cases exhaust the possibilities: if both $\bar\xi\ne0$ and
  $\bar\eta\ne0$, the point is not in $\WF(u_1)\times\WF(u_2)$, so (a) or (b)
  holds; if $\bar\xi=0$ then $\bar\eta\ne0$ and, $\bar x$ being in $\supp
  u_1$, $(\bar y,\bar\eta)\notin\WF(u_2)$, so (b) holds; similarly if
  $\bar\eta=0$.
\end{proof}

\subsection{Restriction}
Write $\R^{p+k}\ni(x,y)$, let $Y=\{y=0\}$ and $\iota:\R^p\to\R^{p+k}$,
$\iota(x)=(x,0)$. We want to define $\iota^*u=u\restriction Y$. No reasonable
such map exists on all distributions (think of $\delta(y)$), but it does under
a condition on the wavefront set: $\WF(u)$ should not meet the \emph{conormal
bundle}
\begin{equation*}
  N^*Y=\{(x,0;0,\eta);\ x\in\R^p,\ \eta\in\R^k\setminus0\}.
\end{equation*}
For $\delta>0$ consider the closed cone
$C_\delta=\{(\xi,\eta)\in\R^{p+k}\setminus0;\ \abs\xi\le\delta\abs\eta\}$ and
let $\mathcal R_\delta$ be the set of $w\in\CmIc(\R^{p+k})$ with $\hat w$
rapidly decreasing on $C_\delta$. If $w\in\mathcal R_\delta$ then, since
$\abs{\hat w}\le C\ang{(\xi,\eta)}^M$,
\begin{equation}\label{microlocal:eq:restriction-int}
  \int\abs{\hat w(\xi,\eta)}\,d\eta\le
  \int_{\abs\eta\le\abs\xi/\delta}C\ang{(\xi,\eta)}^M\,d\eta
  +\int_{\abs\eta>\abs\xi/\delta}C_N\ang\eta^{-N}\,d\eta\le C'\ang\xi^{M+k},
\end{equation}
locally uniformly in $\xi$, so
\begin{equation*}
  g_w(\xi)=\int\hat w(\xi,\eta)\,\dbar\eta\qquad
  (\dbar\eta=(2\pi)^{-k}d\eta)
\end{equation*}
is continuous and polynomially bounded, and we define
$\iota^*w=\cF^{-1}g_w\in\cSp(\R^p)$.

\begin{lemma}\label{microlocal:lem:restriction}
  Let $w\in\mathcal R_\delta$, $\rho\in\CIc(\R^k)$ with $\int\rho=1$, and
  $\rho_\eps(y)=\eps^{-k}\rho(y/\eps)$.
  \begin{enumerate}[label=\textup{(\roman*)}]
  \item For $\phi\in\cS(\R^p)$,
    $\iota^*w(\phi)=\lim_{\eps\to0}w(\phi\boxtimes\rho_\eps)$.
  \item If $w\in\CIc(\R^{p+k})$ then $\iota^*w=w(\cdot,0)$.
  \item $\supp\iota^*w\subset\{x;\ (x,0)\in\supp w\}$; in particular
    $\iota^*w\in\CmIc(\R^p)$, and $\iota^*w=0$ if $\supp w\cap Y=\emptyset$.
  \item For $\phi\in\CIc(\R^p)$, $\phi(x)w\in\mathcal R_{\delta/2}$ and
    $\iota^*(\phi w)=\phi\,\iota^*w$.
  \end{enumerate}
\end{lemma}
\begin{proof}
  (i) By \eqref{microlocal:eq:parseval} in $\R^{p+k}$,
  $w(\phi\boxtimes\rho_\eps)=\iint\hat w(\xi,\eta)\hat\phi(-\xi)
  \hat\rho(-\eps\eta)\,\dbar\xi\,\dbar\eta$. As $\eps\to0$,
  $\hat\rho(-\eps\eta)\to\hat\rho(0)=1$ with $\abs{\hat\rho}\le\norm\rho_{L^1}$,
  and $\abs{\hat w(\xi,\eta)\hat\phi(-\xi)}$ is integrable by
  \eqref{microlocal:eq:restriction-int}; by dominated convergence the limit is
  $\int g_w(\xi)\hat\phi(-\xi)\,\dbar\xi=\iota^*w(\phi)$.

  (ii) is clear from (i), and (iii) follows from (i) too: if
  $(\bar x,0)\notin\supp w$ and $\phi$ is supported near $\bar x$, then
  $\phi\boxtimes\rho_\eps$ vanishes near $\supp w$ for small $\eps$.

  (iv) By \cref{microlocal:lem:fourier-global}(ii),
  $\Sigma(w)$ does not meet the interior of $C_\delta$, which contains
  $C_{\delta/2}$; as $\Sigma(\phi w)\subset\Sigma(w)$,
  \cref{microlocal:lem:fourier-global}(i) shows $\phi w\in\mathcal
  R_{\delta/2}$. The identity follows from (i), since
  $(\phi w)(f\boxtimes\rho_\eps)=w(\phi f\boxtimes\rho_\eps)$.
\end{proof}

\begin{theorem}[Restriction]\label{microlocal:thm:restriction}
  Let $u\in\CmIc(\R^{p+k})$ satisfy $\WF(u)\cap N^*Y=\emptyset$. There are
  $r,\delta>0$ such that $\chi u\in\mathcal R_\delta$ for every
  $\chi=\chi(y)\in\CIc(\R^k)$ supported in $\{\abs y<r\}$, and
  \begin{equation*}
    \iota^*u=\iota^*(\chi u)\in\CmIc(\R^p),\qquad\chi=1\text{ near }0,
  \end{equation*}
  is independent of the choice of $\chi$. The map $u\mapsto\iota^*u$ is
  linear, coincides with $u\mapsto u(\cdot,0)$ on $\CIc(\R^{p+k})$, satisfies
  $\iota^*u(\phi)=\lim_{\eps\to0}u(\phi\boxtimes\rho_\eps)$,
  $\iota^*(\phi u)=\phi\,\iota^*u$ for $\phi\in\CIc(\R^p)$, and
  \begin{equation*}
    \WF(\iota^*u)\subset\iota^*\WF(u)=
    \{(x,\xi)\in\Tz{\R^p};\ (x,0,\xi,\eta)\in\WF(u)\text{ for some }\eta\in\R^k\}.
  \end{equation*}
\end{theorem}
\begin{proof}
  The compact sets
  $Z_\delta=\WF(u)\cap\{\abs y\le\delta\}\cap(\R^{p+k}\times C_\delta)\cap
  \{\abs{(\xi,\eta)}=1\}$ decrease as $\delta\downarrow0$, and their
  intersection is contained in $\WF(u)\cap N^*Y=\emptyset$. Hence
  $Z_\delta=\emptyset$ for some $\delta>0$; take $r=\delta$. If $\supp
  \chi\subset\{\abs y<r\}$ then $\WF(\chi u)\subset\WF(u)\cap\{\abs y<r\}$, so
  $\Sigma(\chi u)\cap C_\delta=\emptyset$ and $\chi u\in\mathcal R_\delta$ by
  \cref{microlocal:lem:fourier-global}(i). If $\chi,\chi'$ are two choices,
  $(\chi-\chi')u\in\mathcal R_\delta$ vanishes near $Y$, so
  $\iota^*(\chi u)=\iota^*(\chi'u)$ by \cref{microlocal:lem:restriction}(iii).
  Linearity is clear, and the other properties, except the last, follow from
  \cref{microlocal:lem:restriction}, since $\chi(y)\rho_\eps(y)=\rho_\eps(y)$
  for small $\eps$. If $u\in\mathcal R_{\delta'}$ for some $\delta'>0$ then
  $\iota^*u$ as defined before \cref{microlocal:lem:restriction} agrees with
  $\iota^*(\chi u)$, by the limit formula.

  For the wavefront set, let $(\bar x,\bar\xi)\in\Tz{\R^p}$ with $(\bar
  x,0,\bar\xi,\eta)\notin\WF(u)$ for all $\eta$. For $\theta\in(0,1)$ let
  \begin{equation*}
    C'_\theta=C_\theta\cup\{(\xi,\eta);\ \xi\in\overline{V_\theta(\bar\xi)}\},
  \end{equation*}
  a closed cone in $\R^{p+k}\setminus0$. As $\theta\downarrow0$ the compact
  sets $C'_\theta\cap\Sphere^{p+k-1}$ decrease to the set of unit vectors
  $(\xi,\eta)$ with $\xi=0$ or $\xi\in\R_+\bar\xi$, which is disjoint from the
  closed cone $\{(\xi,\eta);\ (\bar x,0,\xi,\eta)\in\WF(u)\}$ by hypothesis
  and since $\WF(u)\cap N^*Y=\emptyset$. Hence for some $\theta$ the latter
  cone lies in the open cone $\R^{p+k}\setminus(C'_\theta\cup\{0\})$, and by
  \cref{microlocal:lem:fourier-global}(iii) there are $\phi\in\CIc(\R^p)$,
  $\phi(\bar x)\ne0$, and $\chi(y)$ as above (with small supports) such that
  $\Sigma(\phi\chi u)\cap C'_\theta=\emptyset$. Then
  $\widehat{\phi\chi u}$ is rapidly decreasing on $C'_\theta$; in particular
  $\phi\chi u\in\mathcal R_\theta$, and for $\xi\in V_\theta(\bar\xi)$, where
  $(\xi,\eta)\in C'_\theta$ for all $\eta$,
  \begin{equation*}
    \abs{\widehat{\phi\,\iota^*u}(\xi)}=\Bigl|\int\widehat{\phi\chi u}(\xi,\eta)\,
    \dbar\eta\Bigr|\le C_N\int\ang{(\xi,\eta)}^{-N-k-1}\,d\eta\le C'_N\ang\xi^{-N}.
  \end{equation*}
  Here we used $\phi\iota^*u=\iota^*(\phi\chi u)$, by the independence of
  $\chi$ and \cref{microlocal:lem:restriction}(iv). By
  \cref{microlocal:prop:fourier-char}, $(\bar x,\bar\xi)\notin\WF(\iota^*u)$.
\end{proof}

One can picture this as follows: of $\WF(u)$ only the part over $Y$
survives, and the covector $(\xi,\eta)$ is projected to its component $\xi$
along $Y$; the hypothesis guarantees that no non-zero covector is projected
to zero.

\subsection{Projection}
Let $\pi_1:\R^{p+k}\to\R^p$, $\pi_1(x,y)=x$. For $u\in\CmIc(\R^{p+k})$ define
the push-forward, or ``integral over the fibres'', $(\pi_1)_*u\in\CmIc(\R^p)$
by
\begin{equation*}
  \Bigl(\int u(x,y)\,dy\Bigr)(\phi)=(\pi_1)_*u(\phi)=u(\pi_1^*\phi),\qquad
  \pi_1^*\phi(x,y)=\phi(x),
\end{equation*}
which makes sense because $u$ acts on $\CI(\R^{p+k})$. Thus projection is
dual to pull-back of functions. By duality,
$\supp(\pi_1)_*u\subset\pi_1(\supp u)$, since $\supp\pi_1^*\phi=\pi_1^{-1}(\supp\phi)$.
Moreover $(\pi_1)_*$ maps $\CIc(\R^{p+k})$ to $\CIc(\R^p)$, and
\begin{equation}\label{microlocal:eq:projection-ft}
  \widehat{(\pi_1)_*u}(\xi)=u(e^{-ix\cdot\xi})=\hat u(\xi,0),\qquad
  \phi\,(\pi_1)_*u=(\pi_1)_*\bigl((\pi_1^*\phi)u\bigr).
\end{equation}

\begin{proposition}[Projection]\label{microlocal:prop:projection}
  For $u\in\CmIc(\R^{p+k})$,
  \begin{equation*}
    \WF\bigl((\pi_1)_*u\bigr)\subset\{(x,\xi)\in\Tz{\R^p};\
    (x,y,\xi,0)\in\WF(u)\text{ for some }y\in\R^k\}.
  \end{equation*}
\end{proposition}
\begin{proof}
  Let $(\bar x,\bar\xi)$ be outside the right side. We claim there are
  $r,\theta>0$ such that $\WF(u)$ does not meet
  $\{\abs{x-\bar x}<r\}\times\R^k\times C''_\theta$, where
  $C''_\theta=\{(\xi,\eta);\ \xi\in\overline{V_\theta(\bar\xi)},\
  \abs\eta\le\theta\abs\xi\}$. Otherwise there are
  $(x_j,y_j,\xi_j,\eta_j)\in\WF(u)$ with $\abs{(\xi_j,\eta_j)}=1$,
  $x_j\to\bar x$ and $(\xi_j,\eta_j)\in C''_{1/j}$; since $y_j$ lies in the
  compact set $\pi_2(\supp u)$, a subsequence converges to a point
  $(\bar x,y,\bar\xi/\abs{\bar\xi},0)$ of the closed set $\WF(u)$, a
  contradiction. Now let $\phi\in\CIc(\R^p)$ be supported in $\{\abs{x-\bar
  x}<r\}$ with $\phi(\bar x)\ne0$. Then
  $\Sigma((\pi_1^*\phi)u)\cap C''_\theta=\emptyset$, so by
  \cref{microlocal:lem:fourier-global}(i) and
  \eqref{microlocal:eq:projection-ft},
  $\widehat{\phi(\pi_1)_*u}(\xi)=\widehat{(\pi_1^*\phi)u}(\xi,0)$ is rapidly
  decreasing on $V_\theta(\bar\xi)$. By \cref{microlocal:prop:fourier-char},
  $(\bar x,\bar\xi)\notin\WF((\pi_1)_*u)$.
\end{proof}

\begin{corollary}\label{microlocal:cor:projection-smooth}
  If $u\in\CmIc(\R^{p+k})$ and $\WF(u)\cap\{(x,y,\xi,0);\ \xi\ne0\}=\emptyset$
  then $(\pi_1)_*u\in\CIc(\R^p)$.
\end{corollary}

The covectors $(x,y,\xi,0)$ are exactly those which are conormal to the
fibres $\{x\}\times\R^k$ over which we integrate. Thus the only singularities
which survive integration are those conormal to the fibres; those which are
even partially in the direction of integration are wiped out. This fits with
the fact that integration in \emph{all} variables leaves no singularities
at all.

% ======================================================================
\section{Products of distributions}\label{microlocal:sec:products}

We express the product as ``$u_1(x)u_2(x+y)$ restricted to $y=0$''. Let
$L:\R^{2n}\to\R^{2n}$ be the linear map $L(x,y)=(x,x+y)$, with $\det L=1$
and transpose $DL^t(\xi,\eta)=(\xi+\eta,\eta)$.

\begin{theorem}[Products]\label{microlocal:thm:products}
  Let $u_1,u_2\in\CmIc(\R^n)$ satisfy
  \begin{equation*}
    (x,\xi)\in\WF(u_1)\Longrightarrow(x,-\xi)\notin\WF(u_2).
  \end{equation*}
  Then $w=L^*(u_1\boxtimes u_2)\in\CmIc(\R^{2n})$ satisfies
  $\WF(w)\cap N^*\{y=0\}=\emptyset$, and the product
  \begin{equation*}
    u_1u_2=\iota^*L^*(u_1\boxtimes u_2)\in\CmIc(\R^n)
  \end{equation*}
  has the following properties.
  \begin{enumerate}[label=\textup{(\roman*)}]
  \item $\supp(u_1u_2)\subset\supp u_1\cap\supp u_2$, and for
    $\phi\in\cS(\R^n)$ and $\rho$, $\rho_\eps$ as in
    \cref{microlocal:lem:restriction} (now on $\R^n$), with
    $\check\rho_\eps(t)=\rho_\eps(-t)$,
    \begin{equation*}
      (u_1u_2)(\phi)=\lim_{\eps\to0}u_1\bigl(\phi\cdot(u_2*\check\rho_\eps)\bigr).
    \end{equation*}
  \item If $u_1$ or $u_2$ is smooth, $u_1u_2$ is the usual product.
  \item $(\phi u_1)u_2=\phi(u_1u_2)$ for $\phi\in\CI(\R^n)$.
  \item The wavefront set satisfies
    \begin{align*}
      \WF(u_1u_2)\subset{}&\{(x,\xi);\ x\in\supp u_1,\ (x,\xi)\in\WF(u_2)\}\\
      &\cup\{(x,\xi);\ x\in\supp u_2,\ (x,\xi)\in\WF(u_1)\}\\
      &\cup\{(x,\xi_1+\xi_2);\ (x,\xi_1)\in\WF(u_1),\ (x,\xi_2)\in\WF(u_2)\}.
    \end{align*}
  \end{enumerate}
\end{theorem}
Note that $\xi_1+\xi_2\ne0$ in the last set, by the hypothesis.
\begin{proof}
  By \cref{microlocal:thm:wf-diffeo} (for the linear map $L$),
  \begin{equation*}
    \WF(w)=\{(x,y;\xi+\eta,\eta);\ (x,x+y;\xi,\eta)\in\WF(u_1\boxtimes u_2)\}.
  \end{equation*}
  A point of $N^*\{y=0\}$ has $y=0$ and $\xi+\eta=0$, $\eta\ne0$; by
  \cref{microlocal:prop:exterior}, $(x,x;-\eta,\eta)$ with both components
  non-zero can only lie in $\WF(u_1\boxtimes u_2)$ if $(x,-\eta)\in\WF(u_1)$
  and $(x,\eta)\in\WF(u_2)$, which is excluded. So
  \cref{microlocal:thm:restriction} applies. The support statement follows
  from that theorem, since $(x,0)\in\supp w$ means $(x,x)\in\supp u_1\times\supp
  u_2$. For the formula in (i), by \cref{microlocal:thm:restriction} and the
  definition of $L^*$,
  \begin{equation*}
    (u_1u_2)(\phi)=\lim_{\eps\to0}(u_1\boxtimes u_2)\bigl(\phi(x)\rho_\eps(z-x)\bigr)
    =\lim_{\eps\to0}u_1\bigl(\phi(x)\,u_2(\rho_\eps(\cdot-x))\bigr),
  \end{equation*}
  and $u_2(\rho_\eps(\cdot-x))=(u_2*\check\rho_\eps)(x)$.

  (ii) If $u_2$ is smooth, $u_2*\check\rho_\eps\to u_2$ in $\CI$, so (i) gives
  $u_1(\phi u_2)$. If $u_1$ is smooth,
  $u_1(\phi\cdot(u_2*\check\rho_\eps))=u_2((u_1\phi)*\rho_\eps)\to u_2(u_1\phi)$.

  (iii) is immediate from (i).

  (iv) By \cref{microlocal:thm:restriction} and the formula for $\WF(w)$,
  \begin{equation*}
    \WF(u_1u_2)\subset\{(x,\xi+\eta);\ (x,x;\xi,\eta)\in\WF(u_1\boxtimes u_2)\},
  \end{equation*}
  and \cref{microlocal:prop:exterior} gives the three sets.
\end{proof}

\begin{remark}\label{microlocal:rem:product-fourier}
  If the stronger condition $\Sigma(u_1)\cap(-\Sigma(u_2))=\emptyset$ holds,
  for example if $u_1$ and $u_2$ have small supports near a point $\bar x$ at
  which the hypothesis of the theorem holds
  (\cref{microlocal:lem:fourier-global}(iii)), then the familiar formula
  \begin{equation*}
    \widehat{u_1u_2}(\xi)=\int\hat u_1(\xi-\eta)\hat u_2(\eta)\,\dbar\eta
  \end{equation*}
  holds, with an absolutely convergent integral. Indeed
  $\hat w(\xi,\eta)=\hat u_1(\xi-\eta)\hat u_2(\eta)$, and $w\in\mathcal
  R_\delta$ for small $\delta$: choose open cones $V_i\supset\Sigma(u_i)$ with
  $\overline{V_1}\cap(-\overline{V_2})=\emptyset$; if $\abs\xi\le\delta\abs\eta$
  then either $\eta\notin V_2$, where $\hat u_2(\eta)$ is rapidly decreasing
  and $\abs{\hat u_1(\xi-\eta)}\le C\ang\eta^M$, or $\eta\in V_2$, and then
  $\xi-\eta$, whose direction differs from that of $-\eta$ by at most
  $2\delta$, lies in a closed cone disjoint from $\Sigma(u_1)$ if $\delta$ is
  small, so $\hat u_1(\xi-\eta)$ is rapidly decreasing in $\ang\eta$. The
  formula is then the definition of $\iota^*w$.
\end{remark}

% ======================================================================
\section{Pull-back}\label{microlocal:sec:pullback}

Now let $F:\Omega_1\to\Omega_2$ be a smooth map between open sets
$\Omega_1\subset\R^n$ and $\Omega_2\subset\R^m$; the dimensions may differ.
When can we define $F^*u$ for $u\in\CmIc(\Omega_2)$, and what is its
wavefront set? For a general $F$ there is no consistent definition for all
distributions (consider $F$ constant and $u=\delta$). We assume that $F$ is
\emph{proper}, i.e.\ $F^{-1}(K)$ is compact for every compact $K\subset\Omega_2$,
mostly for convenience: then $F^*\phi=\phi\circ F\in\CIc(\Omega_1)$ for
$\phi\in\CIc(\Omega_2)$. The differential of $F$ defines the cotangent maps
\begin{equation*}
  F^*_x:\R^m=T^*_{F(x)}\Omega_2\to T^*_x\Omega_1=\R^n,\qquad
  \eta\mapsto DF(x)^t\eta,\quad\text{i.e. }
  \xi_i=\sum_j\frac{\pa F_j}{\pa x_i}(x)\eta_j,
\end{equation*}
and we set
\begin{align*}
  N_F&=\{(F(x),\eta)\in\Omega_2\times(\R^m\setminus0);\ x\in\Omega_1,\ DF(x)^t\eta=0\},\\
  F^*W&=\{(x,DF(x)^t\eta);\ (F(x),\eta)\in W\}\quad\text{for }W\subset\Tz{\Omega_2}.
\end{align*}
Note that $F^*$ is a relation rather than a map between cotangent spaces.
$N_F$ is the set of covectors that are mapped to zero; it is empty if $F$ is a
submersion, and for $\iota$ of the previous section it is $N^*Y$.

The construction decomposes $F$ into elementary maps. Let
$U=\{(x,y)\in\Omega_1\times\R^m;\ y+F(x)\in\Omega_2\}$ and
\begin{equation*}
  H:U\to\Omega_1\times\Omega_2,\qquad H(x,y)=(x,y+F(x)),
\end{equation*}
a diffeomorphism with inverse $(x,z)\mapsto(x,z-F(x))$ and $\det DH=1$; it
maps $U\cap\{y=0\}$ onto the graph of $F$. For a function,
$F^*\phi=\iota^*H^*(1\boxtimes\phi)$, $\iota(x)=(x,0)$.

\begin{theorem}[Pull-back]\label{microlocal:thm:pullback}
  Let $F:\Omega_1\to\Omega_2$ be a proper smooth map and let
  $u\in\CmIc(\Omega_2)$ satisfy $\WF(u)\cap N_F=\emptyset$. Let
  $\chi\in\CIc(\Omega_1)$ equal $1$ near the compact set $F^{-1}(\supp u)$.
  Then $H^*(\chi\boxtimes u)\in\CmIc(U)$ satisfies the hypothesis of
  \cref{microlocal:thm:restriction}, and
  \begin{equation*}
    F^*u=\iota^*H^*(\chi\boxtimes u)\in\CmIc(\Omega_1)
  \end{equation*}
  is independent of $\chi$, with $\supp F^*u\subset F^{-1}(\supp u)$. This
  defines a linear map which extends $\phi\mapsto\phi\circ F$ on
  $\CIc(\Omega_2)$, satisfies
  \begin{equation*}
    F^*u(\phi)=\lim_{\eps\to0}u\Bigl(\int\chi(x)\phi(x)\rho_\eps(\cdot-F(x))\,dx\Bigr),
    \qquad\phi\in\cS(\R^n),
  \end{equation*}
  for $\rho\in\CIc(\R^m)$, $\int\rho=1$, agrees with the pull-back of
  \cref{microlocal:lem:diffeo-pullback} if $F$ is a diffeomorphism, and
  \begin{equation*}
    \WF(F^*u)\subset F^*\WF(u)=\bigcup_{x\in\Omega_1}
    \{x\}\times F_x^*\bigl(\WF(u)\cap T^*_{F(x)}\Omega_2\bigr).
  \end{equation*}
\end{theorem}
\begin{proof}
  By \cref{microlocal:prop:exterior}, since $\WF(\chi)=\emptyset$,
  $\WF(\chi\boxtimes u)\subset\{(x,z;0,\eta);\ x\in\supp\chi,\
  (z,\eta)\in\WF(u)\}$. Since $DH(x,y)^t(\xi,\eta)=(\xi+DF(x)^t\eta,\eta)$,
  \cref{microlocal:thm:wf-diffeo} gives
  \begin{equation*}
    \WF\bigl(H^*(\chi\boxtimes u)\bigr)\subset\{(x,y;DF(x)^t\eta,\eta);\
    x\in\supp\chi,\ (y+F(x),\eta)\in\WF(u)\}.
  \end{equation*}
  A point of this set over $y=0$ lies in $N^*\{y=0\}$ only if
  $DF(x)^t\eta=0$ and $(F(x),\eta)\in\WF(u)$, i.e.\ if $(F(x),\eta)\in
  N_F\cap\WF(u)=\emptyset$. Thus \cref{microlocal:thm:restriction} applies,
  and gives $\supp F^*u\subset\{x\in\supp\chi;\ F(x)\in\supp u\}$ and
  $\WF(F^*u)\subset\{(x,DF(x)^t\eta);\ (F(x),\eta)\in\WF(u)\}$.

  If $\chi'$ is another choice, $H^*((\chi-\chi')\boxtimes u)$ is supported in
  $\{(x,y);\ x\in\supp(\chi-\chi'),\ y+F(x)\in\supp u\}$, which does not meet
  $\{y=0\}$ because $\chi=\chi'$ near $F^{-1}(\supp u)$; hence its restriction
  vanishes (\cref{microlocal:lem:restriction}(iii)). By
  \cref{microlocal:thm:restriction}, the formula for $H^*$ in
  \cref{microlocal:lem:diffeo-pullback} (with Jacobian $1$) and the definition
  of the exterior product,
  \begin{equation*}
    F^*u(\phi)=\lim_{\eps\to0}H^*(\chi\boxtimes u)(\phi\boxtimes\rho_\eps)
    =\lim_{\eps\to0}(\chi\boxtimes u)\bigl(\phi(x)\rho_\eps(z-F(x))\bigr),
  \end{equation*}
  which is the stated formula. If $u$ is smooth, the expression under the
  limit is $\int\chi(x)\phi(x)(u*\check\rho_\eps)(F(x))\,dx$, with
  $\check\rho_\eps(t)=\rho_\eps(-t)$, which converges to
  $\int\chi\phi\,(u\circ F)\,dx=\int\phi\,(u\circ F)\,dx$. If $F$ is a diffeomorphism with
  inverse $G$, substituting $x=G(z')$ shows that the function
  $\int\chi\phi(x)\rho_\eps(z-F(x))\,dx$ is
  $\rho_\eps*f$ with $f=\abs{\det DG}\,(\chi\phi)\circ G\in\CIc(\Omega_2)$, which
  converges to $f$ in $\CI$; so $F^*u(\phi)=u(f)$, which is the definition in
  \cref{microlocal:lem:diffeo-pullback} since $\chi\circ G=1$ near $\supp
  u$.
\end{proof}

The hypothesis $\WF(u)\cap N_F=\emptyset$ is natural: without it the right
side of the wavefront set bound would contain zero covectors.
Restriction to a submanifold, products ($u_1u_2=\Delta^*(u_1\boxtimes u_2)$ with
$\Delta(x)=(x,x)$, for which $N_\Delta=\{(x,x;\xi,-\xi)\}$) and
diffeomorphisms are all special cases; see
\cref{microlocal:ex:pullback-special}.

% ======================================================================
\section{Push-forward}\label{microlocal:sec:pushforward}

The dual operation is push-forward. Pull-back maps $\CIc$ to $\CIc$ (for
proper $F$) but is not always defined on distributions. Dually:

\begin{theorem}[Push-forward]\label{microlocal:thm:pushforward}
  Let $F:\Omega_1\to\Omega_2$ be a smooth map from an open subset of $\R^n$ to
  an open subset of $\R^m$. For $u\in\CmIc(\Omega_1)$,
  \begin{equation*}
    F_*u(\phi)=u(F^*\phi),\qquad\phi\in\CI(\Omega_2),
  \end{equation*}
  defines $F_*u\in\CmIc(\Omega_2)$, with $\supp F_*u\subset F(\supp u)$, and
  \begin{multline*}
    \WF(F_*u)\subset\{(F(x),\eta);\ x\in\supp u,\ \eta\ne0,\ DF(x)^t\eta=0\}\\
    \cup\{(F(x),\eta);\ \eta\ne0,\ (x,DF(x)^t\eta)\in\WF(u)\}.
  \end{multline*}
\end{theorem}

Thus $F_*$ is always defined, but even if $u\in\CIc(\Omega_1)$ in general
$F_*u\notin\CIc(\Omega_2)$; if $F$ is a submersion the first set is empty and
$F_*$ maps $\CIc(\Omega_1)$ into $\CIc(\Omega_2)$. The two sets in the bound
are really the same, if one regards $(x,0)$ as belonging to the ``wavefront
set'' of $u$ at every point of $\supp u$.

\begin{proof}
  $F_*u$ is a distribution since $\phi\mapsto\psi F^*\phi$, $\psi\in\CIc(\Omega_1)$
  equal to $1$ near $\supp u$, is continuous; the support statement is clear.
  Formally,
  \begin{equation*}
    F_*u(\phi)=\int u(x)\phi(F(x))\,dx=\iint u(x)\delta(y-F(x))\phi(y)\,dy\,dx,
  \end{equation*}
  so we write $F_*u=(\pi_2)_*\tilde H^*(u\boxtimes\delta)$, where
  $\delta\in\CmIc(\R^m)$, $\tilde H:\Omega_1\times\R^m\to\Omega_1\times\R^m$,
  $\tilde H(x,y)=(x,y-F(x))$, is a diffeomorphism with $\det D\tilde H=1$, and
  $\pi_2(x,y)=y$. To check this identity, note that
  $\tilde H^*(u\boxtimes\delta)\in\CmIc(\Omega_1\times\R^m)$ is supported in the
  compact graph $\{(x,F(x));\ x\in\supp u\}$ and, by
  \cref{microlocal:lem:diffeo-pullback},
  \begin{equation*}
    (\pi_2)_*\tilde H^*(u\boxtimes\delta)(\phi)
    =(u\boxtimes\delta)\bigl(\psi(x)\phi(z+F(x))\bigr)=u\bigl(\psi\,\phi\circ F\bigr)
    =F_*u(\phi).
  \end{equation*}
  Now we estimate the wavefront set using the elementary operations. Since
  $\WF(\delta)=\{0\}\times(\R^m\setminus0)$ (\cref{microlocal:ex:delta}),
  \cref{microlocal:prop:exterior} gives
  \begin{equation*}
    \WF(u\boxtimes\delta)\subset\{(x,0;\xi,\eta);\ (\xi,\eta)\ne0,\
    (x,\xi)\in\WF(u)\text{ or }(\xi=0,\ x\in\supp u)\}.
  \end{equation*}
  Since $D\tilde H(x,y)^t(\xi,\eta)=(\xi-DF(x)^t\eta,\eta)$,
  \cref{microlocal:thm:wf-diffeo} gives
  \begin{equation*}
    \WF\bigl(\tilde H^*(u\boxtimes\delta)\bigr)\subset\{(x,F(x);\xi-DF(x)^t\eta,\eta);\
    (x,\xi)\in\WF(u)\text{ or }(\xi=0,\ x\in\supp u)\}.
  \end{equation*}
  Finally, by \cref{microlocal:prop:projection} (with the roles of the
  factors exchanged), $(y,\eta)\in\WF(F_*u)$ only if $\eta\ne0$ and
  $(x,y;0,\eta)$ lies in the last set for some $x$, i.e.\ $y=F(x)$ and
  $\xi=DF(x)^t\eta$ with either $(x,DF(x)^t\eta)\in\WF(u)$, or
  $DF(x)^t\eta=0$ and $x\in\supp u$.
\end{proof}

The operator $F_*$ has Schwartz kernel $K(y,x)=\delta(y-F(x))$:
$F_*u(y)=\int K(y,x)u(x)\,dx$. Locally $K$ is the pull-back of
$\delta\boxtimes\chi$, $\chi\in\CIc(\Omega_1)$, under the diffeomorphism
$(y,x)\mapsto(y-F(x),x)$, so by \cref{microlocal:prop:exterior} and
\cref{microlocal:thm:wf-diffeo},
\begin{equation*}
  \WF(K)\subset\{(F(x),x;\eta,-DF(x)^t\eta);\ \eta\ne0\}.
\end{equation*}
In the notation of the next section (with the roles of the two factors
interchanged), $\WFp(K)$ is contained in the relation
$\{((F(x),\eta),(x,DF(x)^t\eta))\}$, and the bound of
\cref{microlocal:thm:pushforward} reads
$\WF(F_*u)\subset\WFp(K)\circ\WF(u)\cup\{(y,\eta);\ (y,x;\eta,0)\in\WF(K),\ x\in\supp u\}$.

% ======================================================================
\section{The wavefront relation}\label{microlocal:sec:relation}

One serious application of these results is to general operators with
distributional kernels. Let $\Omega_1\subset\R^n$, $\Omega_2\subset\R^m$ be open
and $K\in\CmIc(\Omega_1\times\Omega_2)$. Then $K$ defines an operator
\begin{equation*}
  A:\CIc(\Omega_2)\to\CmIc(\Omega_1),\qquad Au(\phi)=K(\phi\boxtimes u),
\end{equation*}
formally $Au(x)=\int K(x,y)u(y)\,dy$. We set
\begin{align*}
  \WFp(K)&=\{((x,\xi),(y,\eta));\ (x,y;\xi,-\eta)\in\WF(K)\}
  \subset(\Omega_1\times\R^n)\times(\Omega_2\times\R^m),\\
  \WF_{\Omega_1}(K)&=\{(x,\xi)\in\Tz{\Omega_1};\ (x,y;\xi,0)\in\WF(K)\text{ for some }y\},\\
  \WFp_{\Omega_2}(K)&=\{(y,\eta)\in\Tz{\Omega_2};\ (x,y;0,-\eta)\in\WF(K)\text{ for some }x\}.
\end{align*}
We regard $\WFp(K)$ as a relation: for $W\subset\Tz{\Omega_2}$,
\begin{equation*}
  \WFp(K)\circ W=\{(x,\xi)\in\Tz{\Omega_1};\ ((x,\xi),(y,\eta))\in\WFp(K)
  \text{ for some }(y,\eta)\in W\}.
\end{equation*}

\begin{theorem}[Wavefront relation]\label{microlocal:thm:relation}
  Let $K\in\CmIc(\Omega_1\times\Omega_2)$ and let $\chi\in\CIc(\Omega_1)$
  equal $1$ near the projection of $\supp K$. For every $u\in\CmIc(\Omega_2)$
  with
  \begin{equation}\label{microlocal:eq:relation-cond}
    \WF(u)\cap\WFp_{\Omega_2}(K)=\emptyset
  \end{equation}
  the product $K\cdot(\chi\boxtimes u)$ is defined, and
  \begin{equation*}
    Au=(\pi_1)_*\bigl(K\cdot(\chi\boxtimes u)\bigr)\in\CmIc(\Omega_1)
  \end{equation*}
  extends the operator $A$ above. It satisfies
  \begin{equation*}
    \WF(Au)\subset\WFp(K)\circ\WF(u)\ \cup\ \{(x,\xi);\ (x,y;\xi,0)\in\WF(K)
    \text{ for some }y\in\supp u\}.
  \end{equation*}
\end{theorem}
\begin{proof}
  By \cref{microlocal:prop:exterior},
  $\WF(\chi\boxtimes u)\subset\{(x,y;0,\eta);\ (y,\eta)\in\WF(u)\}$. The
  product condition of \cref{microlocal:thm:products} fails only at points
  $(x,y;0,-\eta)\in\WF(K)$ with $(y,\eta)\in\WF(u)$, which are excluded by
  \eqref{microlocal:eq:relation-cond}. If $u$ is smooth the product is
  $K\cdot(\chi\otimes u)$ and
  $(\pi_1)_*(K\cdot(\chi\otimes u))(\phi)=K(\chi\phi\boxtimes u)=K(\phi\boxtimes
  u)$, so $Au$ agrees with the previous definition. By
  \cref{microlocal:thm:products}(iv), a point of
  $\WF(K\cdot(\chi\boxtimes u))$ has the form $(x,y;\zeta_1+\zeta_2)$ where
  either $\zeta_1=(\xi_1,\eta_1)$ with $(x,y;\zeta_1)\in\WF(K)$, or
  $\zeta_1=0$ and $(x,y)\in\supp K$; and either $\zeta_2=(0,\eta_2)$ with
  $(y,\eta_2)\in\WF(u)$, or $\zeta_2=0$ and $y\in\supp u$. By
  \cref{microlocal:prop:projection}, $(x,\xi)\in\WF(Au)$ requires such a
  point with $\zeta_1+\zeta_2=(\xi,0)$, $\xi\ne0$. If $\zeta_1=0$ the first
  component of $\zeta_1+\zeta_2$ vanishes, which is excluded. If $\zeta_2=0$
  then $(x,y;\xi,0)\in\WF(K)$ with $y\in\supp u$. Otherwise $\xi_1=\xi$ and
  $\eta_1=-\eta_2$, so $(x,y;\xi,-\eta_2)\in\WF(K)$ with
  $(y,\eta_2)\in\WF(u)$, that is $(x,\xi)\in\WFp(K)\circ\WF(u)$.
\end{proof}

\begin{remark}\label{microlocal:rem:relation-psido}
  If $K=K_B$ for some $B\in\Psido m(\R^n)$ with kernel of compact support in
  $\Omega\times\Omega$, then $Au=Bu$ for every $u\in\CmIc(\Omega)$. Indeed, by
  \cref{microlocal:thm:products}(i) with a product mollifier
  $\rho=\rho^{(1)}\boxtimes\rho^{(2)}$ on $\R^{2n}$, and
  $\phi\in\CIc$, $\psi\in\CIc$ equal to $1$ near $\pi_2(\supp K)$,
  \begin{equation*}
    Au(\phi)=\lim_{\eps\to0}K\bigl((\phi\,(\chi*\check\rho^{(1)}_\eps))\boxtimes
    (\psi\,(u*\check\rho^{(2)}_\eps))\bigr)
    =\lim_{\eps\to0}(Bu_\eps)(\phi)=\lim_{\eps\to0}u_\eps(B^t\phi)=u(B^t\phi),
  \end{equation*}
  where $u_\eps=u*\check\rho^{(2)}_\eps\to u$ in $\cSp$, since
  $\chi*\check\rho^{(1)}_\eps=1$ near the projection of $\supp K$ for small $\eps$.
\end{remark}

\begin{corollary}\label{microlocal:cor:relation}
  Let $K\in\CmIc(\Omega_1\times\Omega_2)$.
  \begin{enumerate}[label=\textup{(\roman*)}]
  \item If $\WF(K)\cap\{(x,y;\xi,0)\}=\emptyset$ then $A$ maps $\CIc(\Omega_2)$
    into $\CIc(\Omega_1)$.
  \item If $\WF(K)\cap\{(x,y;0,\eta)\}=\emptyset$ then $A$ extends to a linear
    map $A:\CmIc(\Omega_2)\to\CmIc(\Omega_1)$ with
    $\WF(Au)\subset\WFp(K)\circ\WF(u)\cup\WF_{\Omega_1}(K)$.
  \end{enumerate}
\end{corollary}
\begin{proof}
  Immediate from \cref{microlocal:thm:relation}: in (i), $\WF(u)=\emptyset$ for
  smooth $u$, and in (ii), condition \eqref{microlocal:eq:relation-cond} is
  void.
\end{proof}

\subsection{Applications}
For a pseudodifferential operator $B\in\Psido m(\R^n)$ with kernel of compact
support, \cref{microlocal:prop:wf-kernel} gives
\begin{equation*}
  \WFp(K_B)=\{((x,\xi),(x,\xi));\ (x,\xi)\in\WFp(B)\},
\end{equation*}
a subset of the identity relation. In particular $\WF(K_B)$ contains no
points of the form $(x,y;\xi,0)$ or $(x,y;0,\eta)$, so both parts of
\cref{microlocal:cor:relation} apply, and by
\cref{microlocal:rem:relation-psido} \cref{microlocal:thm:relation} reduces to
$\WF(Bu)\subset\WFp(B)\cap\WF(u)$; this recovers microlocality
(\cref{microlocal:thm:microlocality}) for such operators. The next
chapters will show that many other natural operators, such as the
solution operators of hyperbolic equations, have kernels whose wavefront
sets are well-behaved relations, and \cref{microlocal:thm:relation} then
describes how they propagate singularities.

\begin{remark}
  The restriction to compactly supported kernels is only for simplicity. If
  $K$ is a distribution on $\Omega_1\times\Omega_2$ with \emph{proper}
  support, i.e.\ both projections $\supp K\to\Omega_i$ are proper, then $A$
  maps $\CIc(\Omega_2)$ to distributions of compact support in $\Omega_1$, and
  the results of this section hold for $u\in\CmIc(\Omega_2)$, since only
  $K$ on the compact set $\supp K\cap(\Omega_1\times\supp u)$ is involved.
\end{remark}

% ======================================================================
\section{Exercises}\label{microlocal:sec:exercises}

\begin{exercise}\label{microlocal:ex:residual-S}
  Is it true that every $E\in\Psido{-\infty}(\R^n)$ defines a map from
  $\cSp(\R^n)$ to $\cS(\R^n)$? (Consider $E=\Op(e)$ with $e\in\cS(\R^n_\xi)$
  and $e(0)\ne0$, acting on the constant function $1$.) Compare with
  \cref{microlocal:lem:residual-compact}.
\end{exercise}

\begin{exercise}
  Show that $\cS(\R^n)$ is dense in $L^2(\R^n)$ by proving that if
  $\phi\in\CIc(\R^n)$ is identically equal to $1$ near the origin then
  \begin{equation*}
    u_k(x)=\phi\Bigl(\frac xk\Bigr)\int e^{ix\cdot\xi}\phi\Bigl(\frac\xi k\Bigr)
    \hat u(\xi)\,\dbar\xi\in\cS(\R^n)\quad\text{for }u\in L^2(\R^n),
  \end{equation*}
  and $u_k\to u$ in $L^2(\R^n)$. Can you see any relation to
  pseudodifferential operators here?
\end{exercise}

\begin{exercise}
  For $k\in\N$ show that
  $u\in\cSp(\R^n)$ can be written $u=\sum_{\abs\alpha\le k}D^\alpha u_\alpha$
  with $u_\alpha\in L^2(\R^n)$ if and only if $\ang D^{-k}u\in L^2(\R^n)$, that
  is, if and only if $u\in H^{-k}(\R^n)$.
\end{exercise}

\begin{exercise}
  Let $a\in\CIb(\R^n)$ be real with $a\ge0$, and
  $A=\sum_{j=1}^nD_{x_j}^2+a(x)$. Show that $Au=0$ has no non-trivial solution
  $u\in L^2(\R^n)$. (Show first that $u\in H^2(\R^n)$, then compute
  $\inner{Au}u$.)
\end{exercise}

\begin{exercise}\label{microlocal:ex:density}
  Show that for any open set $\Omega\subset\R^n$, $\CIc(\Omega)$ is dense in
  $\CmIc(\Omega)$ in the weak topology, i.e.\ every $u\in\CmIc(\Omega)$ is the
  limit of a sequence in $\CIc(\Omega)$ in $\cSp(\R^n)$.
\end{exercise}

\begin{exercise}\label{microlocal:ex:empty-wf}
  Let $h\in\CIc(\R^n)$ be supported in $\{1<\abs\eta<2\}$ and not identically
  zero, and let $a(x,\xi)=h(\xi/\ang x)$. Show that $a\in S^0(\R^n;\R^n)$, that
  $\conesupp(a)=\emptyset$, so $\WFp(\Op(a))=\emptyset$, but that
  $a\notin S^{-1}(\R^n;\R^n)$, so $\Op(a)\notin\Psido{-\infty}(\R^n)$.
\end{exercise}

\begin{exercise}
  Show that for $u\in\cSp(\R^n)$
  \begin{equation*}
    \WF(u)=\bigcap\bigl\{\Char_0(A);\ A\in\Psido0(\R^n),\ Au\in\CI(\R^n)\bigr\},
  \end{equation*}
  so that \cref{microlocal:def:wf-compact} extends verbatim.
\end{exercise}

\begin{exercise}
  Compute the wavefront sets of the following distributions: $\delta\in\cSp(\R^n)$;
  $\abs x\in\cSp(\R^n)$; the characteristic function $\chi_{\mathbb B}$ of the
  unit ball, $\chi_{\mathbb B}(x)=1$ for $\abs x\le1$ and $0$ for $\abs x>1$.
  (For the last, use \cref{microlocal:thm:wf-diffeo} to reduce to the
  Heaviside function of one variable, regarded as a function on $\R^n$.)
\end{exercise}

\begin{exercise}
  Let $\Gamma\subset\Tz{\R^n}$ be a closed cone and define
  \begin{equation*}
    \mathcal C^{-\infty}_{c,\Gamma}(\R^n)=\{u\in\CmIc(\R^n);\
    Au\in\CI(\R^n)\ \forall\,A\in\Psido0(\R^n)\text{ with }\WFp(A)\cap\Gamma=\emptyset\}.
  \end{equation*}
  Show that this is the space of $u\in\CmIc(\R^n)$ with $\WF(u)\subset\Gamma$.
  Describe a complete topology on this space with respect to which
  $\CIc(\R^n)$ is a dense subspace.
\end{exercise}

\begin{exercise}
  Give an alternative proof of \cref{microlocal:prop:parametrix}, not using
  global parametrices, along the following lines. Let $z\notin\Char_m(A)$, $a=\sigma_L(A)$,
  and choose $\eps>0$ so small that $\gamma_\eps=\gamma_{z,\eps}$ is supported
  where $\abs a\ge c\abs\xi^m$. Then $b_0=\gamma_\eps/a\in S^{-m}(\R^n;\R^n)$.
  With $B_0=\Op(b_0)$, show that $G=\Id-AB_0\in\Psido0$ has
  $\sigma_0(G)=[1-\gamma_\eps]$, and that with $C=\Op(\gamma_{z,\eps/4})$,
  $CG\in\Psido{-1}$ and $z\notin\WFp((\Id-C)G)$. Writing
  $\Id-AB_0=CG+R_1$, let $B_1\in\Psido{-1}$ be an asymptotic sum of
  $\sum_{j\ge1}(CG)^j$ and $B=B_0(\Id+B_1)$. Show that $AB=\Id+R_2$ with
  $z\notin\WFp(R_2)$. Construct a left parametrix in the same way, or by
  applying the construction to $A^*$ and taking adjoints.
\end{exercise}

\begin{exercise}\label{microlocal:ex:uniqueness}
  Show that microlocal parametrices are essentially unique: if $B$ and $B'$
  satisfy $z\notin\WFp(AB'-\Id)$ and $z\notin\WFp(BA-\Id)$ then
  $z\notin\WFp(B-B')$.
\end{exercise}

\begin{exercise}
  Given $(\bar x,\bar\xi)\in\Tz{\R^n}$, construct $u\in\CmIc(\R^n)$ with
  \begin{equation*}
    \WF(u)=\{(\bar x,t\bar\xi);\ t>0\}.
  \end{equation*}
  (Hint: take $\phi\in\CIc(\R^n)$ with $\int\phi\ne0$ and
  $u(x)=\sum_{k\ge1}k^{-2}\phi(k(x-\bar x))e^{ik^3x\cdot\bar\xi}$, and use the
  Fourier transform. See \cite[Section~8.1]{Hormander-I} for the construction
  of a distribution with any prescribed closed conic wavefront set.)
\end{exercise}

\begin{exercise}
  Suppose $A\in\Psido m(\R^n)$ has Schwartz kernel of compact support. If
  $u\in\CmIc(\R^n)$ use the elementary operations (and
  \cref{microlocal:prop:wf-kernel}) to investigate under what conditions
  $\kappa(x,y)=K_A(x,y)u(y)$, and then $\gamma=(\pi_1)_*\kappa$, make sense.
  What can you say about $\WF(\gamma)$?
\end{exercise}

\begin{exercise}
  Consider the projection $\pi_1:\R^p\times\R^k\to\R^p$. Show that $(\pi_1)_*$
  can be extended to some distributions which do not have compact support,
  for example to
  $\{u\in\cSp(\R^{p+k});\ \supp u\cap(K\times\R^k)\text{ is compact for each compact }K\subset\R^p\}$,
  with values in distributions on $\R^p$ (that is, continuous linear functionals
  on $\CIc(\R^p)$), and that \cref{microlocal:prop:projection} remains valid.
\end{exercise}

\begin{exercise}\label{microlocal:ex:pairing-product}
  Let $u_1,u_2\in\CmIc(\R^n)$ satisfy \eqref{microlocal:eq:pairing-cond}.
  Show that $P(u_1,u_2)=(u_1u_2)(\psi)$ for any $\psi\in\CIc(\R^n)$ equal to $1$
  near $\supp u_1\cup\supp u_2$. (Use \cref{microlocal:thm:products}(i):
  replace $u_2$ by $u_2*\check\rho_\eps$ in $P$, and use
  \cref{microlocal:lem:transpose}.) Deduce that $u_1u_2=u_2u_1$.
\end{exercise}

\begin{exercise}\label{microlocal:ex:pullback-special}
  Check that \cref{microlocal:thm:pullback} contains
  \cref{microlocal:thm:restriction} (with $F=\iota$) and
  \cref{microlocal:thm:products} (with $F=\Delta$, $\Delta(x)=(x,x)$, applied to
  $u_1\boxtimes u_2$): show that the definitions agree, using the limit
  formulas, and that the wavefront set bounds coincide.
\end{exercise}

\section*{Notes}
The definition of the wavefront set via pseudodifferential operators
follows Melrose's notes, while the Fourier transform characterization of
\cref{microlocal:sec:fourier} is H\"ormander's original definition; see
\cite[Chapter~8]{Hormander-I} for the classical treatment of the calculus of
wavefront sets (pull-back, products, push-forward and the wavefront relation
of kernels), including continuity statements in the spaces
$\mathcal D'_\Gamma$ which we have not discussed, and \cite[Chapter~18]{Hormander-III}
for the essential support of pseudodifferential operators.

Compared with the original notes, the main changes are as follows. Essential
supports are used together with compactly supported cut-offs throughout,
because with uniformly bounded symbols $\WFp(A)=\emptyset$ does not imply
$A\in\Psido{-\infty}(\R^n)$. The operations on wavefront sets are all proved
via the Fourier transform characterization and a single restriction theorem;
products are defined as restrictions of $u_1(x)u_2(x+y)$ to $y=0$ (rather than
through the pairing), and the pull-back by a general map is defined through
restriction to a graph, with a limit formula showing consistency with all
the special cases. The wavefront relation theorem includes the term
$\WF_{\Omega_1}(K)$, which is needed in general.
